\section{Lo que traemos de la Sesión 7}

\emph{Advertencia. Dos correcciones de la sesión pasada.} El morfismo canónico entre el límite de un diagrama y el producto de sus objetos va del límite al producto: el límite es un cono sobre el diagrama sin flechas y se factoriza por el producto. La gavilla $\mathcal O_X$ del pegado se obtiene, sobre cada abierto, como un límite inverso o proyectivo de anillos. Un límite directo es un colímite sobre un conjunto dirigido; el pegado no es de ese tipo.

Para los tallos, en cambio, $\mathcal O_{X,x}=\varinjlim_{U\ni x}\mathcal O_X(U)$ es un límite directo o colímite de las restricciones a vecindades cada vez menores.

La sesión anterior probó el \myref{teo:s07-lema-pegado}. Hoy lo usamos con dos piezas, caso en el que el cociclo se cumple automáticamente.

\begin{defi}[Datos de pegado]
\label{defi:s08-datos-pegado}
Unos datos de pegado para una familia de esquemas $\set{X_i}_{i\in I}$ consisten, para cada par $(i,j)\in I\times I$, en un subesquema abierto $X_{ij}\subseteq X_i$, con $X_{ii}=X_i$, y un isomorfismo de esquemas $\varphi_{ij}:X_{ij}\to X_{ji}$ que cumplen:
\begin{enumerate}[label=(P\arabic*), leftmargin=2.1em]
\item $\varphi_{ii}=\id_{X_i}$ para todo $i$;
\item $\varphi_{ji}=\varphi_{ij}^{-1}$ para todo $i,j$;
\item para toda terna $i,j,k$, $\varphi_{ij}(X_{ij}\cap X_{ik})=X_{ji}\cap X_{jk}$ y $\varphi_{ik}=\varphi_{jk}\circ\varphi_{ij}$ en $X_{ij}\cap X_{ik}$.
\end{enumerate}
\end{defi}

Con dos piezas, $I=\set{1,2}$, basta dar abiertos $X_{12}\subseteq X_1$, $X_{21}\subseteq X_2$ y un isomorfismo $\varphi_{12}:X_{12}\to X_{21}$. Se pone $\varphi_{11}=\varphi_{22}=\id$ y $\varphi_{21}=\varphi_{12}^{-1}$: (P1) y (P2) valen por construcción y (P3) se sigue de ellas. \emph{[Hartshorne, II.2, ejercicio 2.12; \cref{obs:s07-dos-piezas-no-cociclo-verificar}.]}

\begin{teo}[Lema de Pegado]
\label{teo:s08-lema-pegado}
Dados datos de pegado para $\set{X_i}_{i\in I}$, existe un esquema $X$ con inmersiones abiertas $\psi_i:X_i\to X$ tales que:
\begin{enumerate}[label=(\alph*)]
\item $X=\bigcup_{i\in I}\psi_i(X_i)$;
\item $\psi_i(X_i)\cap\psi_j(X_j)=\psi_i(X_{ij})$ para todo $i,j$;
\item $\psi_j\circ\varphi_{ij}=\psi_i$ en $X_{ij}$ para todo $i,j$.
\end{enumerate}
Si $(X',\psi_i')$ es otro esquema con inmersiones abiertas que cumplen estas propiedades, existe un único isomorfismo $h:X\to X'$ con $h\circ\psi_i=\psi_i'$ para todo $i$.
\end{teo}
La demostración está en el \cref{teo:s07-lema-pegado}. \emph{[Hartshorne, II.2, ejercicio 2.12; EGA I, \S2.3.]}

\begin{obs}[Morfismos desde un pegado]
\label{obs:s08-morfismos-desde-pegado}
Dar un morfismo $f:X\to Y$ equivale a dar morfismos $f_i:X_i\to Y$ tales que $f_i|_{X_{ij}}=f_j\circ\varphi_{ij}$ para todo $i,j$, por la \cref{prop:s07-pegado-coigualador}. Hoy usaremos esta propiedad con $Y=\Spec A$.
\end{obs}

% Presentación S08, diapositiva 6: coigualador y morfismos desde el pegado.
\begin{center}
\begin{tikzcd}[column sep=huge]
\displaystyle\coprod_{i,j}X_{ij}
  \arrow[r,shift left=.65ex,"\text{inclusión}"]
  \arrow[r,shift right=.65ex,swap,"\varphi_{ij}\text{ e inclusión}"]
& \displaystyle\coprod_iX_i \arrow[r,"\psi"] & X \arrow[r,"f"] & Y
\end{tikzcd}
\end{center}

\section{El plano afín menos el origen}

\begin{ejem}[El plano afín menos el origen]
\label{ejem:s08-plano-menos-origen}
Sea $k$ un campo, $A=k[x,y]$ y $\mathbb A_k^2=\Spec A$. El punto cerrado $(x,y)$ es el único primo que contiene a $(x,y)$. Su complemento
\[
U=\mathbb A_k^2\setminus\set{(x,y)}=D(x)\cup D(y)
\]
es abierto, pues $V(x)\cap V(y)=V(x,y)=\set{(x,y)}$. Con su estructura de subesquema abierto, $U$ es un esquema. Vamos a calcular $\Gamma(U,\mathcal O_U)$. \emph{[Vakil, \S4.4.1 a \S4.4.3.]}
\end{ejem}

% Presentación S08, diapositiva 8: los cuatro dibujos de abiertos del plano.
\begin{center}
\begin{tikzpicture}[every node/.style={font=\small}]
\foreach \a/\col in {0/blue,3.6/green,7.2/green,10.8/orange}{
  \fill[\col!12] (\a,0) rectangle ++(2,2);
  \draw[grisCIMAT] (\a,0) rectangle ++(2,2);
}
\draw[rojoCIMAT,dashed,thick] (1,0)--(1,2);
\draw[rojoCIMAT,dashed,thick] (3.6,1)--(5.6,1);
\draw[rojoCIMAT,dashed,thick] (8.2,0)--(8.2,2);
\draw[rojoCIMAT,dashed,thick] (7.2,1)--(9.2,1);
\draw[grisCIMAT] (11.8,0)--(11.8,2) (10.8,1)--(12.8,1);
\fill[white] (11.8,1) circle (3pt);
\draw[rojoCIMAT,thick] (11.8,1) circle (3pt);
\node at (1,-.4) {$D(x)$};
\node at (1,-.8) {$x\ne0$};
\node at (4.6,-.4) {$D(y)$};
\node at (4.6,-.8) {$y\ne0$};
\node at (8.2,-.4) {$D(x)\cap D(y)=D(xy)$};
\node at (8.2,-.8) {$xy\ne0$};
\node at (11.8,-.4) {$U=D(x)\cup D(y)$};
\node at (11.8,-.8) {todo salvo el origen};
\end{tikzpicture}
\end{center}

\begin{prop}[El plano menos el origen]
\label{prop:s08-secciones-plano-menos-origen}
La restricción
\[
\Gamma(\mathbb A_k^2,\mathcal O)\longrightarrow\Gamma(U,\mathcal O_U)
\]
es un isomorfismo. En particular, $\Gamma(U,\mathcal O_U)=k[x,y]$.
\end{prop}
\begin{proof}
\emph{Las secciones.} La cubierta $U=D(x)\cup D(y)$ tiene traslape $D(xy)$. Por los axiomas de identidad y pegado, una sección global es un par $(f,g)\in A_x\times A_y$ cuyas imágenes coinciden en $A_{xy}$. Como $A$ es un dominio, las localizaciones son inyectivas y podemos considerar
\[
A\subseteq A_x,A_y\subseteq A_{xy},
\qquad \Gamma(U,\mathcal O_U)=A_x\cap A_y.
\]
Si $h=p/x^a=q/y^b$ pertenece a la intersección, con $p,q\in A$ y $a,b\geq0$, entonces $py^b=qx^a$. El anillo $A$ es un dominio de factorización única, y $x,y$ son primos distintos. De $x^a\mid py^b$ y $x\nmid y$ se obtiene $x^a\mid p$, luego $h\in A$. La inclusión opuesta es inmediata: la restricción es un isomorfismo.

\emph{Advertencia.} La inclusión $A\hookrightarrow A_x$ usa que $A$ es un dominio; para un anillo cualquiera la localización puede perder elementos.
\end{proof}

\begin{obs}[Lectura geométrica]
\label{obs:s08-hartogs}
Los polos de $p/x^a$ están en la recta $x=0$, y los de $q/y^b$ en la recta $y=0$. Si representan la misma función, sus polos tendrían que estar en el origen. En dimensión dos, los polos de una función forman una curva y un punto no lo es. Toda función regular en $\mathbb A_k^2$ menos el origen se extiende al plano: es el análogo algebraico del lema de Hartogs. En $\mathbb A_k^1\setminus\set{0}$ sí hay funciones nuevas, como $1/x$. Para $n\geq2$ se usa la cubierta $D(x_1),\ldots,D(x_n)$ de $\mathbb A_k^n$ menos el origen; véase el \cref{ejer:s08-afin-menos-origen}.
\end{obs}

\emph{Comparación topológica.} El plano real $\mathbb R^2$ es contráctil, mientras que $\mathbb R^2\setminus\set{0}$ tiene el tipo de homotopía de $S^1$ y grupo fundamental $\mathbb Z$. En el caso complejo, $\mathbb C^2\setminus\set{0}$ tiene el tipo de homotopía de $S^3$, que es simplemente conexa; lo distingue de $\mathbb C^2$ su tercer grupo de homología $H_3=\mathbb Z$ (o su $\pi_3$).

\clearpage
\subsection{El morfismo canónico}
\label{sec:s08-morfismo-canonico}

\begin{defi}[El morfismo canónico $\eta_X$]
\label{defi:s08-morfismo-canonico}
Sea $X$ un esquema y $R=\Gamma(X,\mathcal O_X)$. El morfismo canónico $\eta_X:X\to\Spec R$ se define de la siguiente manera. Para $x\in X$, denotamos por $\mathfrak m_x$ el ideal maximal de $\mathcal O_{X,x}$ y por $\kappa(x)=\mathcal O_{X,x}/\mathfrak m_x$ su cuerpo residual.
\begin{enumerate}[label=(\roman*)]
\item En puntos,
\[
\eta_X(x)=\set*{a\in R\given a_x\in\mathfrak m_x}
=\ker\bigl(R\longrightarrow\mathcal O_{X,x}\longrightarrow\kappa(x)\bigr).
\]
Es un primo porque $\kappa(x)$ es un campo.
\item En abiertos, para $f\in R$,
\[
\eta_X^{-1}(D(f))=X_f
:=\set*{x\in X\given f_x\notin\mathfrak m_x},
\]
el abierto donde $f$ no se anula.
\item En gavillas, sobre cada básico,
\[
\eta_f^\sharp:\mathcal O_{\Spec R}(D(f))=R_f
\longrightarrow\mathcal O_X(X_f),
\qquad
\frac{a}{f^n}\longmapsto\frac{a|_{X_f}}{(f|_{X_f})^n}.
\]
\end{enumerate}
El conjunto $X_f$ es abierto porque un germen invertible tiene una inversa en algún entorno. Las inversas locales se pegan, de modo que $f|_{X_f}$ es invertible y la fórmula tiene sentido. En tallos, $a/s\mapsto a_x/s_x$ es un homomorfismo local. \emph{[Hartshorne, II, ejercicio 2.4; EGA I, 1.6.3.]}
\end{defi}

\begin{ejer}[Verificación del morfismo canónico]
\label{ejer:s08-verificacion-morfismo-canonico}
Sea $X$ un esquema y $R=\Gamma(X,\mathcal O_X)$. Muestra que $\eta_X:X\to\Spec R$ es un morfismo de esquemas verificando:
\begin{enumerate}[label=(\roman*)]
\item $\eta_X(x)$ es un ideal primo de $R$.
\item $\eta_X$ es continuo, pues $\eta_X^{-1}(D(f))=X_f$ es abierto.
\item $\eta_X^\sharp$ está bien definido: $f|_{X_f}$ es invertible en $\mathcal O_X(X_f)$ y los homomorfismos $R_f\to\mathcal O_X(X_f)$ son compatibles con las restricciones.
\item En tallos, $(\eta_X)_x^\sharp:R_{\eta_X(x)}\to\mathcal O_{X,x}$ es un homomorfismo local.
\end{enumerate}
\end{ejer}
% S08b, diapositivas 3--5: esta verificación se desarrolla en clase.
\begin{proof}[Desarrollo en clase]
\emph{Puntos y topología.} El ideal $\eta_X(x)$ es primo como núcleo de la evaluación en el cuerpo $\kappa(x)$; es propio porque $1$ se envía a $1\ne0$. Si $x\in X_f$, el germen $f_x$ es unidad. Elige un representante local $g$ de su inversa. La igualdad $g_xf_x=1$ vale en un entorno $V$ de $x$, pues dos secciones con el mismo germen coinciden cerca de $x$. Así $V\subseteq X_f$, luego $X_f$ es abierto. La igualdad $\eta_X^{-1}(D(f))=X_f$ y la base de abiertos $D(f)$ prueban la continuidad.

\emph{Gavillas.} Las inversas locales de $f$ en $X_f$ son únicas y por ello se pegan en una inversa de $f|_{X_f}$. La propiedad universal de la localización factoriza la restricción de manera única:
% S08b, diapositiva 4: localización y compatibilidad con restricciones.
\begin{center}
\begin{minipage}{.45\textwidth}
\centering
\begin{tikzcd}[column sep=large]
R\arrow[r,"\mathrm{res}"]\arrow[d,swap,"\mathrm{loc}"]
 &\mathcal O_X(X_f)\\
R_f\arrow[ur,dashed,swap,"\eta_f^\sharp"]&
\end{tikzcd}
\end{minipage}\hfill
\begin{minipage}{.51\textwidth}
\centering
\begin{tikzcd}[column sep=large]
R_f\arrow[r,"\eta_f^\sharp"]\arrow[d]
 &\mathcal O_X(X_f)\arrow[d,"\mathrm{res}"]\\
R_g\arrow[r,swap,"\eta_g^\sharp"]&\mathcal O_X(X_g)
\end{tikzcd}
\end{minipage}
\end{center}
Si $D(g)\subseteq D(f)$, también $X_g\subseteq X_f$. Ambos caminos son mapas $R_f\to\mathcal O_X(X_g)$ que coinciden en $R$. La unicidad de la localización da la igualdad. Por el \cref{teo:s05-toda-gavilla-sobre-base-extiende}, los mapas sobre los básicos definen un morfismo de gavillas. La compatibilidad con las restricciones es indispensable.

\emph{Tallos.} Sea $\mathfrak p=\eta_X(x)$. Si $s\notin\mathfrak p$, el germen $s_x$ es unidad; por tanto el mapa de gérmenes se factoriza por $R_{\mathfrak p}$:
\nopagebreak
% S08b, diapositiva 5: gérmenes, localización y cuerpo residual.
\begin{center}
\begin{tikzcd}[column sep=large]
R\arrow[r,"\text{germen}"]\arrow[d,swap,"\mathrm{loc}"]
 &\mathcal O_{X,x}\arrow[r]&\kappa(x)\\
R_{\mathfrak p}\arrow[ur,dashed,swap,"(\eta_X)_x^\sharp"]&&
\end{tikzcd}
\end{center}
El mapa inducido es
\[
(\eta_X)_x^\sharp:R_{\mathfrak p}\longrightarrow\mathcal O_{X,x},
\qquad \frac{a}{s}\longmapsto\frac{a_x}{s_x}.
\]
Como $s_x$ es unidad, la imagen de $a/s$ está en $\mathfrak m_x$ si y solo si $a_x\in\mathfrak m_x$, es decir, si y solo si $a\in\mathfrak p$. La preimagen de $\mathfrak m_x$ es $\mathfrak pR_{\mathfrak p}$: el homomorfismo es local. Así $\eta_X$ es un morfismo de espacios localmente anillados y, entre esquemas, un morfismo de esquemas.
\end{proof}

En una carta afín $V\cong\Spec B$, el morfismo es $\Spec(\rho_V)$ para la restricción
\[
\rho_V:R\longrightarrow\Gamma(V,\mathcal O_X)\cong B.
\]
En puntos manda $\mathfrak q$ a $\rho_V^{-1}(\mathfrak q)$ y en básicos induce las localizaciones de $\rho_V$, como en el \cref{teo:s06-todo-morfismo-anillos-b-induce}. Esta descripción muestra también que la construcción es independiente de la cubierta.

\begin{prop}[Qué dice el morfismo canónico]
\label{prop:s08-propiedades-morfismo-canonico}
Sea $X$ un esquema y $R=\Gamma(X,\mathcal O_X)$.
\begin{enumerate}[label=(\alph*)]
\item $\eta_X$ es un morfismo de espacios localmente anillados que en secciones globales es la identidad de $R$.
\item Para todo anillo $B$ y todo $\varphi:B\to R$ hay un único $g:X\to\Spec B$ que induce $\varphi$ en secciones globales, a saber $g=\Spec(\varphi)\circ\eta_X$. Es decir,
\[
\Hom(X,\Spec B)\cong\Hom_{\mathbf{Ring}}(B,\Gamma(X,\mathcal O_X)).
\]
\item Si $h:X\to Y$ es un morfismo, entonces
\[
\eta_Y\circ h=\Spec(\Gamma(h))\circ\eta_X.
\]
\item Si $\psi:X\to Y$ es un isomorfismo, $\eta_X$ es un isomorfismo si y solo si $\eta_Y$ lo es.
\item $X$ es afín si y solo si $\eta_X$ es un isomorfismo.
\end{enumerate}
\emph{[Hartshorne, II, ejercicio 2.4; EGA I, 1.6.3.]}
\end{prop}
\begin{proof}
\emph{La identidad en secciones globales.} La verificación está en el \cref{ejer:s08-verificacion-morfismo-canonico}. Con $f=1$, se tiene $X_1=X$ y el mapa $R\to\mathcal O_X(X)$ es la identidad; esto prueba (a).

\emph{Naturalidad y propiedad universal.} Para (c), sea $h:X\to Y$, con $h(x)=y$. Como el homomorfismo $h_x^\sharp:\mathcal O_{Y,y}\to\mathcal O_{X,x}$ es local, para $a\in\Gamma(Y,\mathcal O_Y)$,
\[
\Gamma(h)(a)\in\eta_X(x)
\iff h_x^\sharp(a_y)\in\mathfrak m_x
\iff a_y\in\mathfrak m_y
\iff a\in\eta_Y(y).
\]
Así $\Gamma(h)^{-1}(\eta_X(x))=\eta_Y(h(x))$, que da la igualdad en puntos. En gavillas, sobre $D(a)$ ambos caminos envían una fracción $b/a^n$ a la misma fracción de secciones restringidas, por la compatibilidad de $h^\sharp$ con restricciones y gérmenes. Luego los mapas de gavillas también coinciden.
% S08b, diapositiva 8: naturalidad general y con destino afín.
\begin{center}
\begin{minipage}{.48\textwidth}
\centering
\begin{tikzcd}[column sep=large]
X\arrow[r,"\eta_X"]\arrow[d,swap,"h"]
 &\Spec\Gamma(X,\mathcal O_X)\arrow[d,"\Spec\Gamma(h)"]\\
Y\arrow[r,swap,"\eta_Y"]&\Spec\Gamma(Y,\mathcal O_Y)
\end{tikzcd}
\end{minipage}\hfill
\begin{minipage}{.48\textwidth}
\centering
\begin{tikzcd}[column sep=large]
X\arrow[r,"\eta_X"]\arrow[d,swap,"g"]
 &\Spec\Gamma(X,\mathcal O_X)\arrow[d,"\Spec\Gamma(g)"]\\
\Spec B\arrow[r,swap,"\id"]&\Spec B
\end{tikzcd}
\end{minipage}
\end{center}
Aquí hemos identificado $\Gamma(\Spec B,\mathcal O)$ con $B$. En efecto, para $X=\Spec B$, la fórmula en puntos es
\[
\eta_{\Spec B}(\mathfrak p)
=\set*{b\in B\given b/1\in\mathfrak pB_{\mathfrak p}}
=\mathfrak p,
\]
y en gavillas $B_f\to B_f$ es la identidad. Por tanto $\eta_{\Spec B}=\id$.

Dado $\varphi:B\to R$, el compuesto $g=\Spec(\varphi)\circ\eta_X$ induce $\varphi$, por (a). Recíprocamente, la naturalidad aplicada a cualquier $g:X\to\Spec B$ da $g=\Spec(\Gamma(g))\circ\eta_X$. Por tanto el homomorfismo en secciones globales determina $g$, lo que prueba (b).
% S08b, diapositiva 7: factorización universal y correspondencia de Hom.
\begin{center}
\begin{minipage}{.47\textwidth}
\centering
\begin{tikzcd}[column sep=large,row sep=large]
X\arrow[r,"\eta_X"]\arrow[dr,swap,"g"]
 &\Spec R\arrow[d,dashed,"\exists!\,\Spec(\varphi)"]\\
&\Spec B
\end{tikzcd}
\end{minipage}\hfill
\begin{minipage}{.49\textwidth}
\centering
\begin{tikzcd}[row sep=large]
\Hom(X,\Spec B)\arrow[d,shift left=.8ex,"\Gamma"]\\
\Hom_{\mathbf{Ring}}(B,R)
 \arrow[u,shift left=.8ex,"\varphi\mapsto\Spec(\varphi)\circ\eta_X"]
\end{tikzcd}
\end{minipage}
\end{center}
Todo morfismo de $X$ a un esquema afín se factoriza de manera única por $\eta_X$. Es la unidad de la adjunción entre $\Gamma$ y $\Spec$: la correspondencia anterior envía $g$ a $\Gamma(g)$ y $\varphi$ a $\Spec(\varphi)\circ\eta_X$.

\emph{Isomorfismos y afinidad.} Si $\psi:X\to Y$ es un isomorfismo, también lo son $\Gamma(\psi)$ y $\Spec(\Gamma(\psi))$. El cuadrado de naturalidad tiene flechas verticales invertibles:
% S08b, diapositiva 9: transporte del morfismo canónico por un isomorfismo.
\begin{center}
\begin{tikzcd}[column sep=large]
X\arrow[r,"\eta_X"]\arrow[d,swap,"\psi","\sim"']
 &\Spec\Gamma(X,\mathcal O_X)
 \arrow[d,"\Spec\Gamma(\psi)","\sim"']\\
Y\arrow[r,swap,"\eta_Y"]&\Spec\Gamma(Y,\mathcal O_Y)
\end{tikzcd}
\end{center}
Por ello
\[
\eta_X=\Spec(\Gamma(\psi))^{-1}\circ\eta_Y\circ\psi,
\]
y $\eta_X$ es un isomorfismo si y solo si lo es $\eta_Y$, que es (d). En puntos, si $Y=\Spec B$ y $\psi(x)=\mathfrak p$, para $a=\Gamma(\psi)(b)$ se tiene $a_x=\psi_x^\sharp(b/1)$. Este isomorfismo local lleva $b/1$ al maximal si y solo si $b\in\mathfrak p$. Así $\eta_X(x)$ es el ideal transportado $\Gamma(\psi)(\mathfrak p)$.

Si $X$ es afín, elige $\psi:X\xrightarrow{\sim}\Spec B$. Como $\eta_{\Spec B}=\id$, (d) muestra que $\eta_X$ es un isomorfismo. A la inversa, si $\eta_X$ es un isomorfismo, identifica $X$ con el esquema afín $\Spec\Gamma(X,\mathcal O_X)$.

La presentación da también el argumento mediante Yoneda. Para todo esquema $Z$, conmuta el triángulo
% S08b, diapositiva 10: argumento de adjunción y Yoneda.
\begin{center}
\begin{tikzcd}[column sep=large,row sep=large]
\Hom(Z,X)\arrow[r,"\eta_X\circ-"]\arrow[dr,swap,"\Gamma"]
 &\Hom(Z,\Spec\Gamma(X,\mathcal O_X))
 \arrow[d,"\sim"]\\
&\Hom_{\mathbf{Ring}}(\Gamma(X,\mathcal O_X),\Gamma(Z,\mathcal O_Z)).
\end{tikzcd}
\end{center}
La flecha vertical es biyectiva por (b). Si $X$ es afín, la diagonal también lo es, aplicando (b) tras identificar $X$ con un espectro. Por tanto $\eta_X\circ-$ induce biyecciones en todos los conjuntos $\Hom(Z,-)$, y Yoneda implica que $\eta_X$ es un isomorfismo.
\end{proof}

\subsection{Aplicación al plano menos el origen}

\begin{obs}[El caso afín: identidad e inclusión]
\label{obs:s08-canonico-identidad-inclusion}
Para $X=\Spec A$, el morfismo canónico es la identidad bajo $\Gamma(X,\mathcal O_X)=A$. Si $U\subseteq\Spec A$ es abierto, con inclusión $j$ y restricción $\rho:A\to\Gamma(U,\mathcal O_U)$, la naturalidad da
\[
j=\Spec(\rho)\circ\eta_U.
\]
\end{obs}
% S08b, diapositiva 11: distinguir el canónico y la inclusión.
\begin{center}
\begin{tikzcd}[column sep=large,row sep=large]
U\arrow[r,"\eta_U"]\arrow[d,swap,"j"]
 &\Spec\Gamma(U,\mathcal O_U)\arrow[d,"\Spec(\rho)"]\\
\Spec A\arrow[r,swap,"\eta_{\Spec A}=\id"]&\Spec A
\end{tikzcd}
\end{center}
En puntos, para $\mathfrak p\in U$,
\[
\Spec(\rho)(\eta_U(\mathfrak p))
=\rho^{-1}\set*{g\in\Gamma(U,\mathcal O_U)\given g_{\mathfrak p}\in\mathfrak m_{\mathfrak p}}
=\mathfrak p.
\]
El germen de $a|_U$ en $\mathfrak p$ es el de $a$. Si $\rho$ es un isomorfismo, $\eta_U$ se identifica con la inclusión $j$ al identificar sus destinos mediante $\Spec(\rho)$. Antes de hacerlo, los destinos son distintos.

\begin{coro}
\label{coro:s08-plano-menos-origen-no-afin}
El abierto $U=\mathbb A_k^2\setminus\set{(x,y)}$ no es afín.
\end{coro}
\begin{proof}
Por la \cref{prop:s08-secciones-plano-menos-origen}, la restricción $\rho:k[x,y]\to\Gamma(U,\mathcal O_U)$ es un isomorfismo. La \cref{obs:s08-canonico-identidad-inclusion} da el cuadrado
% S08b, diapositiva 12: el canónico del plano perforado y la inclusión.
\begin{center}
\begin{tikzcd}[column sep=large,row sep=large]
U\arrow[r,"\eta_U"]\arrow[d,swap,"j"]
 &\Spec\Gamma(U,\mathcal O_U)\arrow[d,"\Spec(\rho)","\sim"']\\
\mathbb A_k^2\arrow[r,swap,"\eta_{\mathbb A_k^2}=\id"]&\mathbb A_k^2
\end{tikzcd}
\end{center}
Si $U$ fuera afín, $\eta_U$ sería un isomorfismo por la \cref{prop:s08-propiedades-morfismo-canonico}(e). Como $\Spec(\rho)$ lo es, también lo sería $j=\Spec(\rho)\circ\eta_U$. Pero $j$ omite el origen $(x,y)$ y no es sobreyectiva.

En la identificación $\Gamma(U,\mathcal O_U)=A=k[x,y]$, la fórmula es
\[
\eta_U(\mathfrak p)
=\set*{a\in A\given a_{\mathfrak p}\in\mathfrak pA_{\mathfrak p}}
=\mathfrak p,
\]
pues $a_{\mathfrak p}\in\mathfrak pA_{\mathfrak p}$ equivale a $a\in\mathfrak p$. Las funciones globales de $U$ son las del plano entero, cuyo espectro recupera el origen ausente.
\end{proof}

\section{La recta proyectiva por cartas}

Con la técnica anterior, calculemos secciones globales de un esquema construido por pegado. Primero construiremos sus cartas y sus transiciones.

\begin{defi}[La recta proyectiva $\mathbb P_A^1$]
\label{defi:s08-recta-proyectiva}
Sea $A$ un anillo. Sean $U_0=\Spec A[t]$ y $U_1=\Spec A[s]$, dos copias de la recta afín $\mathbb A_A^1$, llamadas cartas. Los abiertos $U_{01}=D(t)\subseteq U_0$ y $U_{10}=D(s)\subseteq U_1$ tienen anillos de funciones $A[t,t^{-1}]$ y $A[s,s^{-1}]$. El isomorfismo de anillos
\[
A[s,s^{-1}]\longrightarrow A[t,t^{-1}],\qquad s\longmapsto t^{-1},
\]
define un isomorfismo de esquemas $\varphi_{01}:U_{01}\to U_{10}$. Por la \cref{defi:s08-datos-pegado}, son datos de pegado. La recta proyectiva $\mathbb P_A^1$ es el esquema del \cref{teo:s08-lema-pegado}, con inmersiones abiertas $\psi_0:U_0\to\mathbb P_A^1$ y $\psi_1:U_1\to\mathbb P_A^1$. \emph{[Vakil, \S4.4.6 y \S4.4.10; Hartshorne, II.2, ejercicio 2.12.]}
\end{defi}

\begin{obs}[Coordenadas]
\label{obs:s08-coordenadas}
Un punto $[x_0:x_1]$ con $x_0\ne0$ tiene coordenada $t=x_1/x_0$; cuando $x_1\ne0$, tiene coordenada $s=x_0/x_1$. En el traslape se cumple $s=t^{-1}$, la transición de la \cref{defi:s08-recta-proyectiva}. Si $A=k$ y $a\in k^\times$, los puntos $t=a$ y $s=a^{-1}$ de ambas cartas se identifican.
\end{obs}

% Presentación S08, diapositiva 13: círculo y las dos proyecciones.
\begin{center}
\begin{minipage}{.47\textwidth}
\centering
\begin{tikzpicture}[scale=1,every node/.style={font=\small}]
\draw[grisCIMAT,thick] (0,0) circle (1.2);
\draw[blue!60!black,thick] (-2.4,-1.2)--(2.4,-1.2) node[right] {$t$};
\draw[green!45!black,thick] (-2.4,1.2)--(2.4,1.2) node[right] {$s$};
\begin{scope}
\clip (-2.4,-1.2) rectangle (2.4,1.2);
\foreach \x/\y in {-.98/.69,.98/.69,-.98/-.69,.98/-.69}{
 \pgfmathsetmacro{\projt}{2.4*\x/(1.2-\y)}
 \pgfmathsetmacro{\projs}{2.4*\x/(1.2+\y)}
 \draw[blue!60!black,dashed] (0,1.2)--(\projt,-1.2);
 \draw[green!45!black,dashed] (0,-1.2)--(\projs,1.2);
}
\end{scope}
\foreach \x/\y in {-.98/.69,.98/.69,-.98/-.69,.98/-.69}{
 \fill[rojoCIMAT] (\x,\y) circle (1.8pt);
}
\fill[rojoCIMAT] (0,1.2) circle (2.5pt);
\fill[rojoCIMAT] (0,-1.2) circle (2.5pt);
\node[above] at (0,1.2) {$s=0$};
\node[below] at (0,-1.2) {$t=0$};
\end{tikzpicture}
\end{minipage}\hfill
\begin{minipage}{.49\textwidth}
\small
La carta $t$ proyecta desde $s=0$ sobre la recta tangente en $t=0$. La carta $s$ proyecta desde $t=0$ sobre la tangente en $s=0$. El traslape es el círculo menos los dos polos, donde $ts=1$. Sobre $\mathbb R$ se ve un círculo; sobre $\mathbb C$, la esfera de Riemann.
\end{minipage}
\end{center}

\begin{obs}[El morfismo estructural]
\label{obs:s08-morfismo-estructural}
Los morfismos $U_0\to\Spec A$ y $U_1\to\Spec A$ inducidos por $A\to A[t]$ y $A\to A[s]$ coinciden en el traslape, pues $s\mapsto t^{-1}$ es un homomorfismo de $A$-álgebras. La \cref{obs:s08-morfismos-desde-pegado} da $\pi:\mathbb P_A^1\to\Spec A$; sus fibras se estudian en el \cref{sec:s08-puntos}.
\end{obs}

% S08b, diapositiva 15: propiedad universal del pegado y triángulo de anillos.
Sean $\iota_0:U_{01}\hookrightarrow U_0$ y $\iota_1:U_{10}\hookrightarrow U_1$. La condición para pegar $f_0:U_0\to Y$ y $f_1:U_1\to Y$ es $f_0\circ\iota_0=f_1\circ\iota_1\circ\varphi_{01}$. La \cref{obs:s08-morfismos-desde-pegado} da entonces un único $f:\mathbb P_A^1\to Y$:
\begin{center}
\begin{minipage}{.59\textwidth}
\centering
\begin{tikzcd}[column sep=large,row sep=large]
U_{01}\arrow[r,"\iota_0"]\arrow[d,swap,"\iota_1\circ\varphi_{01}"]
 &U_0\arrow[d,"\psi_0"]\arrow[ddr,bend left=25,"f_0"]&\\
U_1\arrow[r,swap,"\psi_1"]\arrow[drr,bend right=20,swap,"f_1"]
 &\mathbb P_A^1\arrow[dr,dashed,"\exists!\,f"]&\\
&&Y
\end{tikzcd}
\end{minipage}\hfill
\begin{minipage}{.37\textwidth}
\centering
\begin{tikzcd}[column sep=large,row sep=large]
A\arrow[r]\arrow[dr,swap]
 &A[s,s^{-1}]\arrow[d,"s\mapsto t^{-1}"]\\
&A[t,t^{-1}]
\end{tikzcd}
\end{minipage}
\end{center}
Para $Y=\Spec A$, el triángulo de anillos verifica la compatibilidad de los morfismos estructurales: la transición es un homomorfismo de $A$-álgebras. De este pegado se obtiene $\pi$.

\begin{prop}[$\mathbb P_A^1$ como conjunto]
\label{prop:s08-proyectiva-como-conjunto}
Se tiene $\mathbb P_A^1=\psi_0(U_0)\amalg\psi_1(V(s))$, con
\[
V(s)=\Spec(A[s]/(s))\cong\Spec A.
\]
Así, $\mathbb P_A^1$ es la recta afín $\mathbb A_A^1$ más una copia de $\Spec A$, la sección del infinito. Esta sección es un punto cuando $A$ es un campo.
\end{prop}
\begin{proof}
Por (a) y (b) del \cref{teo:s08-lema-pegado}, el complemento de $\psi_0(U_0)$ es
\[
\psi_1(U_1\setminus D(s))=\psi_1(V(s)).
\]
El cociente $A[s]/(s)\cong A$ identifica este conjunto con $\Spec A$.
\end{proof}

% Presentación S08, diapositiva 14: mismas piezas, dos pegados.
\begin{center}
\begin{tikzpicture}[every node/.style={font=\small}]
\begin{scope}
\node at (0,1.7) {recta con origen doble};
\draw[blue!60!black,thick] (-2.2,1)--(2.2,1) node[right] {$X_1$};
\draw[green!45!black,thick] (-2.2,0)--(2.2,0) node[right] {$X_2$};
\foreach \x in {-1.8,-1.2,-.6,.6,1.2,1.8}{\draw[grisCIMAT,dotted,semithick] (\x,0)--(\x,1);}
\fill[rojoCIMAT] (0,1) circle (2.5pt) node[above] {$O_1$};
\fill[rojoCIMAT] (0,0) circle (2.5pt) node[below] {$O_2$};
\node at (0,-.9) {$\varphi^\sharp(y)=x$};
\end{scope}
\begin{scope}[xshift=7cm]
\node at (0,1.7) {$\mathbb P_k^1$};
\draw[blue!60!black,thick] (-2.2,1)--(2.2,1) node[right] {$t$};
\draw[green!45!black,thick] (-2.2,0)--(2.2,0) node[right] {$s$};
\draw[grisCIMAT,dotted,semithick] (-2,1)--(-.4,0) (-.4,1)--(-2,0) (.4,1)--(2,0) (2,1)--(.4,0);
\fill[rojoCIMAT] (0,1) circle (2.5pt) node[above] {$t=0$};
\fill[rojoCIMAT] (0,0) circle (2.5pt) node[below] {$s=0$};
\node at (0,-.9) {$\varphi^\sharp(y)=x^{-1}$};
\end{scope}
\end{tikzpicture}
\end{center}
Las piezas $\mathbb A_k^1$ y el traslape $D(x)$ son los mismos. Con la identidad, cerca de $O_1$ y de $O_2$ aparecen los mismos puntos. Con la inversión, cerca de $t=0$ en una carta es lejos en la otra. El morfismo de pegado distingue las dos construcciones.
\clearpage
\section{Secciones globales y no afinidad}

\begin{prop}
\label{prop:s08-secciones-proyectiva}
Para todo anillo $A$, se tiene $\Gamma(\mathbb P_A^1,\mathcal O)=A$: el homomorfismo $A\to\Gamma(\mathbb P_A^1,\mathcal O)$ inducido por $\pi$ es un isomorfismo. \emph{[Vakil, \S4.4.8.]}
\end{prop}
\begin{proof}
Una sección global es un par $(f,g)\in A[t]\times A[s]$ con $f=g(t^{-1})$ en $A[t,t^{-1}]$, por la \cref{defi:s07-gavilla-estructural-pegado}. Es el igualador de las dos restricciones:
% S08b, diapositiva 19: las secciones globales como igualador.
\begin{center}
\begin{tikzcd}[column sep=huge]
\Gamma(\mathbb P_A^1,\mathcal O)\arrow[r,hook]
 &A[t]\times A[s]
 \arrow[r,shift left=.8ex,"{(f,g)\mapsto f}"]
 \arrow[r,shift right=.8ex,swap,"{(f,g)\mapsto g(t^{-1})}"]
 &A[t,t^{-1}]
\end{tikzcd}
\end{center}
Escribamos
\[
f=\sum_{n\geq0}a_nt^n,\qquad g(t^{-1})=\sum_{n\geq0}b_nt^{-n}.
\]
El anillo de Laurent es libre como $A$-módulo con base $\set{t^n}_{n\in\mathbb Z}$. Comparar coeficientes da $a_n=b_n=0$ para $n>0$ y $a_0=b_0$. Luego $f=g=a_0\in A$. Cada $a\in A$ da, a su vez, la sección compatible $(a,a)$.

\emph{Advertencia.} $A[t]\hookrightarrow A[t,t^{-1}]$ es inyectivo para cualquier $A$, sin suponer que sea un dominio.
\end{proof}

\begin{coro}
\label{coro:s08-proyectiva-no-afin}
Si $A\ne0$, la recta proyectiva $\mathbb P_A^1$ no es afín.
\end{coro}
\begin{proof}
Sea $R=\Gamma(\mathbb P_A^1,\mathcal O)$ y $\varphi:A\to R$ el isomorfismo de la \cref{prop:s08-secciones-proyectiva}. Por la \cref{prop:s08-propiedades-morfismo-canonico}(b), aplicada a $g=\pi$, se tiene
\[
\pi=\Spec(\varphi)\circ\eta_{\mathbb P_A^1}.
\]
% S08b, diapositiva 21: identificar pi con el morfismo canónico.
\begin{center}
\begin{tikzcd}[column sep=large,row sep=large]
\mathbb P_A^1\arrow[r,"\eta_{\mathbb P_A^1}"]\arrow[dr,swap,"\pi"]
 &\Spec R\arrow[d,"\Spec(\varphi)","\sim"']\\
&\Spec A
\end{tikzcd}
\end{center}
Si $\mathbb P_A^1$ fuera afín, su morfismo canónico sería un isomorfismo por la \cref{prop:s08-propiedades-morfismo-canonico}(e). Como $\Spec(\varphi)$ también lo es, $\pi$ sería un isomorfismo y, en particular, biyectivo.

Como $A\ne0$, existe un primo $\mathfrak r\subseteq A$. En la carta $U_0$ consideremos
\[
P_0=\mathfrak r A[t]+(t),\qquad
P_1=\mathfrak r A[t]+(t-1).
\]
Ambos son primos de la misma carta $U_0$, pues $A[t]/P_i\cong A/\mathfrak r$. Son distintos: la imagen de $t-1$ en $A[t]/P_0$ es $-1\ne0$. La inmersión abierta $\psi_0$ es inyectiva, de modo que también son distintos en $\mathbb P_A^1$; no hace falta cambiar de carta. Además, $P_i\cap A=\mathfrak r$. En $U_0$, el morfismo $\pi$ es $\Spec$ de la inclusión $A\to A[t]$, por lo que envía cada primo a su contracción. Así,
\[
\pi(P_0)=\mathfrak r=\pi(P_1),
\]
y $\pi$ no es inyectivo. En términos de gérmenes, $a\in A$ pertenece al maximal del tallo en $P_i$ si y solo si $a\in\mathfrak r$. Esto confirma la fórmula del morfismo canónico.
\end{proof}

% Presentación S08, diapositiva 18: no sobreyectividad y no inyectividad.
\begin{center}
\begin{tikzpicture}[every node/.style={font=\small}]
\fill[orange!12] (0,0) rectangle (1.6,1.6);
\draw[grisCIMAT] (0,0) rectangle (1.6,1.6);
\draw[rojoCIMAT,thick,fill=white] (.8,.8) circle (2.5pt);
\draw[->,blue!60!black,thick] (1.95,.8)--(3.15,.8) node[midway,above] {$\eta_U$};
\fill[orange!12] (3.5,0) rectangle (5.1,1.6);
\draw[grisCIMAT] (3.5,0) rectangle (5.1,1.6);
\fill[rojoCIMAT] (4.3,.8) circle (2.5pt) node[above right,xshift=-2.5pt,yshift=-2.5pt] {$\scriptstyle(x,y)$};
\node at (.8,-.4) {$U$};
\node at (4.3,-.4) {$\mathbb A_k^2=\Spec\Gamma(U,\mathcal O_U)$};
\node[align=center] at (2.55,-1.25) {\footnotesize Falta el origen: el mapa no es sobreyectivo.};
\begin{scope}[xshift=11cm]
\draw[green!45!black,thick] (0,.8)--(0,2.3);
\node[left] at (0,1.6) {$\mathbb P_A^1$};
\foreach \y/\lab in {2.1/P_0,1.6/P_1,1.1/\infty}{\fill[rojoCIMAT] (0,\y) circle (2.5pt) node[right] {$\lab$};}
\draw[->,blue!60!black,thick] (0,.7)--(0,.05) node[midway,left] {$\pi$};
\draw[grisCIMAT] (-1.4,-.2)--(1.4,-.2) node[right] {$\Spec A$};
\fill[rojoCIMAT] (0,-.2) circle (2.5pt) node[below] {$\mathfrak r$};
\node[align=center] at (0,-1.25) {\footnotesize $P_0\ne P_1$ tienen misma imagen: el mapa no es inyectivo.};
\end{scope}
\end{tikzpicture}
\end{center}
Un esquema afín se reconstruye con sus funciones. En $U$ las funciones no detectan el hueco del origen; en $\mathbb P_A^1$ no distinguen los dos puntos elegidos.

\begin{obs}
\label{obs:s08-constantes-compacidad}
Es el análogo de que las únicas funciones holomorfas en $\mathbb{CP}^1$ son constantes; en el caso complejo, esto se explica por la compacidad. Una función entera no constante tiene en $\infty$ un polo si es un polinomio y una singularidad esencial si no lo es, como $e^z$. Si fuera holomorfa también en $\infty$, sería constante por Liouville. La versión algebraica de la compacidad es la propiedad, que se estudiará en la Sesión 15. La recta con origen doble tiene secciones globales $k[x]$, que la distinguen de $\mathbb P_k^1$; véase el \cref{sec:s08-dos-rectas}. \emph{[Vakil, ejercicio 4.4.B.]}
\end{obs}

\section{Los puntos de la recta proyectiva}
\label{sec:s08-puntos}

\begin{obs}[Qué buscamos]
\label{obs:s08-objetivos-puntos}
Describiremos los puntos de la recta proyectiva sobre un cuerpo, sobre $\mathbb Z$ y sobre un anillo de valoración discreta. Las cartas reducen el problema a clasificar primos de anillos de polinomios. Después estudiaremos la reducción de puntos racionales a una fibra especial, construiremos $\mathbb P_A^n$ y clasificaremos los pegados de dos rectas.
\end{obs}

\subsection{Lo que se usa hoy}

\subsubsection{Definiciones y resultados previos}
Recordemos las \cref{defi:s01-espectro-primo,defi:s02-s-definimos-v-s-p,defi:s02-f-abierto-distinguido-asociado-f}: los puntos de $\Spec A$ son sus ideales primos, los cerrados son $V(I)=\set{\mathfrak p\in\Spec A\given I\subseteq\mathfrak p}$ y los abiertos básicos son $D(f)=\set{\mathfrak p\given f\notin\mathfrak p}$. Por el \cref{coro:s05-gavilla-estructural}, $\mathcal O(D(f))=A_f$; además, $D(f)\cong\Spec A_f$ por el \cref{coro:s03-inclusion-inducida-identifica-f-d}.

Un punto $x$ es cerrado cuando $\set{x}$ es cerrado, y es genérico de un cerrado $Z$ cuando $\overline{\set{x}}=Z$. Por la \cref{prop:s02-todo-p-cerradura-punto-p}, $\overline{\set{\mathfrak p}}=V(\mathfrak p)$: los puntos cerrados del espectro son los ideales máximos y, si $A$ es un dominio, $(0)$ es su punto genérico. La \cref{defi:s02-decimos-q-especializacion-p-q} describe las especializaciones.

\begin{obs}[La especialización es local]
\label{obs:s08-especializacion-local}
Si $y\in\overline{\set{x}}$ y un abierto $U$ contiene a $y$, contiene también a $x$: toda vecindad de $y$ corta a $\set{x}$. Así, $x\notin U$ implica $\overline{\set{x}}\cap U=\varnothing$.
\end{obs}

Por la \cref{obs:s03-cuerpo-residual}, el cuerpo residual de $\mathfrak p$ es
\[
\kappa(\mathfrak p)=A_{\mathfrak p}/\mathfrak pA_{\mathfrak p}=\Frac(A/\mathfrak p).
\]
En un esquema, $\kappa(x)=\mathcal O_{X,x}/\mathfrak m_x$, y una carta afín conserva este cuerpo. Un homomorfismo $\varphi:A\to B$ induce $\Spec\varphi$, dado por $\mathfrak q\mapsto\varphi^{-1}(\mathfrak q)$; véanse el \cref{lema:s01-preimagen-primo} y el \cref{ejer:s02-b-morfismo-anillos-demuestra-1}. Además,
\[
(\Spec\varphi)^{-1}(D(f))=D(\varphi(f)).
\]

Una inmersión abierta identifica su fuente con un subesquema abierto de la meta; en particular, es inyectiva en puntos. Las cartas de la \cref{defi:s08-recta-proyectiva} son inmersiones abiertas. Su morfismo estructural se construyó en la \cref{obs:s08-morfismo-estructural}; para sus secciones y su no afinidad, véanse \cref{prop:s08-secciones-proyectiva,coro:s08-proyectiva-no-afin}. El criterio de afinidad y el plano menos el origen se trataron en \cref{prop:s08-propiedades-morfismo-canonico,obs:s08-canonico-identidad-inclusion,prop:s08-secciones-plano-menos-origen,coro:s08-plano-menos-origen-no-afin}.

\begin{defi}[Esquema sobre $A$, morfismo estructural y fibra]
\label{defi:s08-esquema-sobre-anillo-fibra}
Un esquema sobre $A$ es un esquema $X$ con un morfismo $\pi:X\to\Spec A$, llamado morfismo estructural. Para $\mathbb A_A^1=\Spec A[t]$, la inclusión $A\to A[t]$ induce $\pi(\mathfrak q)=\mathfrak q\cap A$. La fibra sobre $y\in\Spec A$ es el conjunto $\pi^{-1}(y)$. Aquí solo se usa como conjunto de puntos; la fibra como esquema, $X\times_{\Spec A}\Spec\kappa(y)$, se construirá en las Sesiones 12 y 14.
\end{defi}

\begin{defi}[Grado de un punto]
\label{defi:s08-grado-punto}
Sea $k$ un campo y $x$ un punto de un esquema sobre $k$ cuyo cuerpo residual es una extensión finita de $k$. Su grado es $[\kappa(x):k]$, donde la estructura de $k$-álgebra viene del morfismo estructural. Un punto de grado uno se llama racional.
\end{defi}

\begin{defi}[Polinomios mónicos, irreducibles y primitivos]
\label{defi:s08-polinomios-primitivos}
Sea $B$ un dominio. Un polinomio de $B[t]$ es mónico si su coeficiente principal es uno. Es irreducible si no es una unidad y no es producto de dos no unidades. El contenido de un polinomio no nulo de $\mathbb Z[t]$ es el máximo común divisor de sus coeficientes; el polinomio es primitivo si su contenido es uno.
\end{defi}

\subsubsection{Hechos de álgebra}

\begin{prop}[Hecho A, polinomios sobre un campo]
\label{prop:s08-polinomios-campo}
Si $k$ es un campo, $k[t]$ es un dominio de ideales principales. Sus primos son $(0)$ y los $(f)$, con $f$ irreducible; estos últimos son máximos y tienen un único generador mónico. Si $\deg f=d$, el campo $k[t]/(f)$ tiene $k$-base $1,t,\ldots,t^{d-1}$. En cualquier dominio de ideales principales, todo primo no nulo es máximo. \emph{[Lang, Algebra, cap.~IV, \S1; Dummit--Foote, \S8.2.]}
\end{prop}

\begin{prop}[Hecho B, primos de una localización]
\label{prop:s08-primos-localizacion}
Sea $B$ un anillo y $S\subseteq B$ multiplicativo. La aplicación $\mathfrak q\mapsto S^{-1}\mathfrak q$ es una biyección entre los primos de $B$ disjuntos de $S$ y $\Spec S^{-1}B$. Su inversa es la contracción por $B\to S^{-1}B$. Es el \cref{teo:s03-correspondencia-localizaciones}. \emph{[Atiyah--Macdonald, prop.~3.11(iv).]}
\end{prop}

\begin{prop}[Hecho C, primos de un cociente]
\label{prop:s08-primos-cociente}
Sea $B$ un anillo e $I$ un ideal. La aplicación $\mathfrak q\mapsto\mathfrak q/I$ es una biyección entre los primos de $B$ que contienen a $I$ y $\Spec(B/I)$. Su inversa es la contracción por $B\to B/I$; véase el \cref{ejer:s01-i-cociente-demuestra-i-inyectivo}. \emph{[Atiyah--Macdonald, cap.~1.]}
\end{prop}

\begin{prop}[Hecho D, lema de Gauss]
\label{prop:s08-gauss}
El producto de dos polinomios primitivos de $\mathbb Z[t]$ es primitivo. \emph{[Lang, Algebra, cap.~IV, \S2.]} En consecuencia:
\begin{enumerate}[label=(G\arabic*)]
\item Todo irreducible de $\mathbb Q[t]$ de grado positivo es, salvo un factor de $\mathbb Q^\times$, un polinomio primitivo $g\in\mathbb Z[t]$, único salvo signo e irreducible en $\mathbb Z[t]$. Recíprocamente, un primitivo de grado positivo irreducible en $\mathbb Z[t]$ es irreducible en $\mathbb Q[t]$. Como $\mathbb Z[t]$ es un dominio de factorización única, $(g)$ es primo.
\item Si $g\in\mathbb Z[t]$ es primitivo y divide a $h\in\mathbb Z[t]$ en $\mathbb Q[t]$, también lo divide en $\mathbb Z[t]$; equivalentemente,
\[
g\mathbb Q[t]\cap\mathbb Z[t]=g\mathbb Z[t].
\]
\end{enumerate}
\end{prop}

\begin{prop}[Hecho E, cuerpos finitos]
\label{prop:s08-cuerpos-finitos}
Sea $K$ un campo de característica $p>0$ y $q$ una potencia de $p$.
\begin{enumerate}[label=(E\arabic*)]
\item El Frobenius $x\mapsto x^q$ es un homomorfismo de anillos $K\to K$.
\item Si $K$ tiene $q$ elementos, todo $x\in K$ cumple $x^q=x$.
\item Si $K$ es finito, el grupo multiplicativo $K^\times$ es cíclico.
\end{enumerate}
\emph{[Lang, Algebra, cap.~V, \S5; Dummit--Foote, \S14.3.]}
\end{prop}

\begin{lema}[Fibra de la recta afín]
\label{lema:s08-fibra-recta-afin}
Sean $A$ un anillo, $\mathfrak p\in\Spec A$ y $\kappa=\kappa(\mathfrak p)$. La fibra sobre $\mathfrak p$ de la contracción $\Spec A[t]\to\Spec A$ está en biyección con $\Spec\kappa[t]$. La biyección lleva $\mathfrak r\in\Spec\kappa[t]$ a su imagen inversa por $A[t]\to\kappa[t]$. Lo mismo vale con $s$ en lugar de $t$.
\end{lema}
\begin{proof}
Localizamos primero en $S=A\setminus\mathfrak p$. Los primos $\mathfrak q$ de $A[t]$ con $\mathfrak q\cap A=\mathfrak p$ corresponden a los primos de $A_{\mathfrak p}[t]$ que contienen a $\mathfrak pA_{\mathfrak p}[t]$, por la \cref{prop:s08-primos-localizacion}. Al reducir por ese ideal, la \cref{prop:s08-primos-cociente} los identifica con los primos de
\[
A_{\mathfrak p}[t]/\mathfrak pA_{\mathfrak p}[t]=\kappa[t].
\]
Recíprocamente, un primo que contiene a $\mathfrak pA_{\mathfrak p}[t]$ contrae a $\mathfrak pA_{\mathfrak p}$ en $A_{\mathfrak p}$, porque este ideal es máximo. Su contracción a $A$ es $\mathfrak p$. Las dos correspondencias inversas son contracciones y su compuesto es el mapa anunciado.
\end{proof}

\begin{lema}[Cerrados y cartas]
\label{lema:s08-cerrados-cartas}
Si $X$ está cubierto por abiertos $U_i$, un subconjunto $Z$ es cerrado si y solo si $Z\cap U_i$ es cerrado en $U_i$ para todo $i$. En particular, un punto es cerrado si y solo si es cerrado en cada carta que lo contiene.
\end{lema}
\begin{proof}
Una intersección de un cerrado con $U_i$ es cerrada en $U_i$. Recíprocamente, si todas las intersecciones son cerradas, cada $U_i\setminus Z$ es abierto en $U_i$ y en $X$; su unión es $X\setminus Z$. Para un punto, las cartas que no lo contienen tienen intersección vacía con él.
\end{proof}

\subsection{Sobre un cuerpo}

\begin{prop}[Los puntos de $\mathbb A_k^1$]
\label{prop:s08-puntos-recta-afin-campo}
Sea $k$ un campo. Los puntos de $\mathbb A_k^1=\Spec k[t]$ son $(0)$ y los $(f)$ con $f$ mónico irreducible, uno por cada tal polinomio. \emph{[Vakil, \S4.4.6; \cref{prop:s01-recta-afin}.]}
\begin{enumerate}[label=(\alph*)]
\item $(0)$ es el punto genérico, no es cerrado y $\kappa((0))=k(t)$.
\item Cada $(f)$ es cerrado, con $\kappa((f))=k[t]/(f)$ y grado $\deg f$.
\item Los puntos de grado uno son los $(t-a)$, con $a\in k$.
\end{enumerate}
\end{prop}
\begin{proof}
La \cref{prop:s08-polinomios-campo} lista los primos y sus cocientes. Como $k[t]$ es un dominio, la cerradura de $(0)$ es todo el espectro; no es máximo, pues $(0)\subsetneq(t)$. Los demás primos son máximos y sus cocientes ya son campos, de dimensión $\deg f$ sobre $k$. Un mónico de grado uno es precisamente $t-a$.
\end{proof}

\begin{prop}[Los puntos de $\mathbb P_k^1$]
\label{prop:s08-puntos-proyectiva-cuerpo}
\label{ejem:s08-puntos-proyectiva-cuerpo}
Sea $k$ un campo e $\infty=\psi_1((s))$. Por la \cref{prop:s08-proyectiva-como-conjunto},
\[
\mathbb P_k^1=\mathbb A_k^1\amalg\set{\infty}.
\]
Sus puntos son el genérico $(0)$, los puntos cerrados $(f)$ con $f\in k[t]$ mónico irreducible y el punto cerrado $\infty$. Se tiene $\kappa((0))=k(t)$, $\kappa((f))=k[t]/(f)$, de grado $\deg f$, y $\kappa(\infty)=k$, de grado uno. \emph{[Vakil, \S4.4.6 y \S4.4.7.]}
\end{prop}
\begin{proof}
La lista procede de las \cref{prop:s08-proyectiva-como-conjunto,prop:s08-puntos-recta-afin-campo}. El genérico $(0)$ está en $D(t)$ y corresponde al genérico $(0)$ de la segunda carta; su cerradura contiene ambas cartas. No es cerrado, pues ya no lo es en $U_0$.

Para $(f)\ne(t)$, la localización por $t$ da un ideal máximo de $k[t,t^{-1}]$. En la segunda carta corresponde a un irreducible no divisible por $s$: al sustituir $t=s^{-1}$ y limpiar denominadores se obtiene un polinomio de término constante no nulo. Así el punto es cerrado en ambas cartas. El punto $(t)$ no está en la segunda carta y $\infty=(s)$ no está en la primera. El \cref{lema:s08-cerrados-cartas} prueba su cerradura global. Los cuerpos residuales se calculan en las cartas afines.
\end{proof}

\begin{coro}[Los puntos cerrados de $\mathbb P_{\mathbb C}^1$ y de $\mathbb P_{\mathbb R}^1$]
\label{coro:s08-puntos-complejos-reales}
\begin{enumerate}[label=(\alph*)]
\item Sobre $\mathbb C$, los puntos cerrados son los $(t-z)$, con $z\in\mathbb C$, e $\infty$, todos de grado uno. Como conjunto son $\mathbb C\cup\set{\infty}$, la esfera de Riemann.
\item Sobre $\mathbb R$, son los $(t-a)$ con $a\in\mathbb R$ y $\infty$, de grado uno, y los $(t^2+bt+c)$ con $b^2<4c$, de grado dos y cuerpo residual $\mathbb C$. Asociar a cada polinomio sus raíces complejas, y a $\infty$ el conjunto $\set{\infty}$, da una biyección con las órbitas de $\operatorname{Gal}(\mathbb C/\mathbb R)=\set{\id,z\mapsto\bar z}$ en $\mathbb C\cup\set{\infty}$.
\end{enumerate}
\emph{[Vakil, \S4.4.7.]}
\end{coro}
\begin{proof}
Por el teorema fundamental del álgebra, los irreducibles complejos son lineales y los reales son lineales o cuadráticos de discriminante negativo. En este último caso, evaluar en una raíz no real identifica el cociente con $\mathbb C$. Las raíces forman un par conjugado; las raíces reales y el infinito son órbitas de un solo elemento. Aplicamos la \cref{prop:s08-puntos-proyectiva-cuerpo}.
\end{proof}

% Presentación S08, diapositiva 21: órbita de la conjugación.
\begin{center}
\begin{tikzpicture}[every node/.style={font=\small}]
\draw[grisCIMAT] (-2.5,0)--(2.5,0) node[right] {$\mathbb R$};
\draw[grisCIMAT] (0,-1.5)--(0,1.5);
\foreach \x in {-1.8,.7,1.8}{\fill[rojoCIMAT] (\x,0) circle (2pt);}
\fill[rojoCIMAT] (-.8,1.1) circle (2.5pt) node[left] {$z$};
\fill[rojoCIMAT] (-.8,-1.1) circle (2.5pt) node[left] {$\bar z$};
\draw[->,blue!60!black,dashed,thick] (-.8,1.05) to[bend left=25] (-.8,-1.05);
\node[right,blue!60!black] at (.25,.7) {$\operatorname{Gal}(\mathbb C/\mathbb R)$};
\end{tikzpicture}
\end{center}
\emph{Advertencia.} Un punto cerrado de grado dos es el par $\set{z,\bar z}$, no dos puntos. El polinomio $t^2+1$ da un solo punto de $\mathbb P_{\mathbb R}^1$, aunque $i$ y $-i$ son dos números.

\begin{lema}[Factorización de $t^{q^n}-t$]
\label{lema:s08-factorizacion-frobenius}
Sea $q$ una potencia de un primo y $n\geq1$. En $\mathbb F_q[t]$, el polinomio $t^{q^n}-t$ es el producto de todos los mónicos irreducibles cuyo grado divide a $n$, cada uno una sola vez.
\end{lema}
\begin{proof}
La derivada de $P=t^{q^n}-t$ es $-1$, por lo que no hay factores repetidos. Sea $f$ mónico irreducible de grado $d$, $K=\mathbb F_q[t]/(f)$ y $\alpha$ la clase de $t$. Por la \cref{prop:s08-polinomios-campo}, $K$ tiene $q^d$ elementos. El Frobenius $F(x)=x^q$ fija $\mathbb F_q$, y $F^d=\id_K$ por la \cref{prop:s08-cuerpos-finitos}.

Si $d\mid n$, resulta $\alpha^{q^n}=\alpha$ y $f\mid P$. Recíprocamente, si $f\mid P$, entonces $F^n(\alpha)=\alpha$. Como $\alpha$ genera $K$ sobre $\mathbb F_q$, $F^n=\id_K$. La ciclicidad de $K^\times$ implica que $q^d-1$ divide a $q^n-1$. Escribamos $n=ad+r$, con $0\leq r<d$. Módulo $q^d-1$, tenemos $q^n-1\equiv q^r-1$; si $r>0$, este es un entero positivo menor que $q^d-1$, contradicción. Por tanto $r=0$, es decir, $d\mid n$.
\end{proof}

\begin{prop}[Puntos cerrados de $\mathbb P_{\mathbb F_q}^1$]
\label{prop:s08-puntos-cerrados-campo-finito}
Sea $q$ una potencia de un primo, y sea $M_d$ el número de puntos cerrados de $\mathbb P_{\mathbb F_q}^1$ de grado $d$, donde el grado de $x$ es $[\kappa(x):\mathbb F_q]$. Para todo $n\geq1$,
\[
\sum_{d\mid n}dM_d=q^n+1.
\]
En particular, $M_1=q+1$ y $M_2=(q^2-q)/2$.
\end{prop}
\begin{proof}
Sea $N_d$ el número de polinomios mónicos irreducibles de grado $d$ en $\mathbb F_q[t]$. Por el \cref{lema:s08-factorizacion-frobenius}, comparar grados da
\[
\sum_{d\mid n}dN_d=q^n.
\]
Los puntos cerrados de grado $d$ de $\mathbb A_{\mathbb F_q}^1$ son esos $N_d$, por la \cref{prop:s08-puntos-recta-afin-campo}. El infinito añade un punto de grado uno, de modo que $M_1=N_1+1$ y $M_d=N_d$ para $d\geq2$. Así se obtiene la fórmula. Para $n=1$ da $M_1=q+1$, y para $n=2$ da $M_1+2M_2=q^2+1$, de donde resulta el valor de $M_2$. \emph{[Lang, Algebra, cap.~V, \S5.]}
\end{proof}

\begin{obs}[Lectura geométrica]
\label{obs:s08-orbitas-frobenius}
Un punto cerrado de grado $d$ es una órbita de Frobenius de $d$ puntos de $\mathbb P^1(\overline{\mathbb F}_q)$, como $\set{z,\bar z}$ es una órbita de conjugación sobre $\mathbb R$. La órbita cae en $\mathbb P^1(\mathbb F_{q^n})$ cuando $d\mid n$ y aporta $d$ puntos. Como este conjunto tiene $q^n+1$ puntos $[a:b]$, la suma cuenta las órbitas por su tamaño.
\end{obs}

\subsection{Las fibras de \texorpdfstring{$\mathbb P_A^1$}{P1A} sobre un anillo cualquiera}

\begin{prop}[Fibras de $\mathbb P_A^1$]
\label{prop:s08-fibras-proyectiva}
Sean $A$ un anillo, $\mathfrak p\in\Spec A$, $\kappa=\kappa(\mathfrak p)$ y $\pi:\mathbb P_A^1\to\Spec A$ el morfismo estructural. Hay una biyección entre el conjunto $\pi^{-1}(\mathfrak p)$ y el conjunto de puntos de $\mathbb P_\kappa^1$, que lleva:
\begin{itemize}
\item el primo $\mathfrak pA[t]$ de $U_0$ al genérico $(0)$ de $\mathbb P_\kappa^1$;
\item cada primo de $U_0$ sobre $\mathfrak p$ correspondiente a $(\bar f)\subseteq\kappa[t]$, con $\bar f$ mónico irreducible, a ese punto $(\bar f)$;
\item el único punto $\infty_{\mathfrak p}$ de $\psi_1(V(s))$ sobre $\mathfrak p$ al punto $\infty$.
\end{itemize}
\end{prop}
\begin{proof}
En $U_0$, la fibra se identifica con $\Spec\kappa[t]$ por el \cref{lema:s08-fibra-recta-afin}. El núcleo de $A[t]\to\kappa[t]$ es $\mathfrak pA[t]$, de modo que este primo corresponde a $(0)$. Por la \cref{prop:s08-proyectiva-como-conjunto}, el complemento de $U_0$ es una copia de $\Spec A$ y aporta exactamente $\infty_{\mathfrak p}$. Este corresponde en $U_1$ a $(\mathfrak p,s)$. Las identificaciones de las dos cartas respetan $s=t^{-1}$, por lo que recuperan la lista de la \cref{prop:s08-puntos-proyectiva-cuerpo} sobre $\kappa$. Solo afirmamos aquí una biyección de conjuntos; la construcción de la fibra como esquema queda para más adelante.
\end{proof}

\subsection{Sobre \texorpdfstring{$\mathbb Z$}{Z}}

\begin{prop}[{Los primos de $\mathbb Z[t]$}]
\label{prop:s08-primos-z-t}
Los ideales primos de $\mathbb Z[t]$ son exactamente los siguientes:
\begin{enumerate}[label=(\alph*)]
\item $(0)$;
\item $(g)$, con $g\in\mathbb Z[t]$ primitivo, irreducible y de grado positivo, único salvo signo;
\item $(p)$, con $p$ un número primo;
\item $(p,g)$, con $p$ primo y $g\in\mathbb Z[t]$ mónico cuya reducción $\bar g\in\mathbb F_p[t]$ es irreducible.
\end{enumerate}
Los dos primeros tipos contraen a $(0)$ en $\mathbb Z$ y los dos últimos a $(p)$. En el último, el ideal depende solo de $\bar g$, no del levantamiento mónico $g$. \emph{[Vakil, ejercicio 3.2.Q; Sesión 1.]}
\end{prop}
\begin{proof}
Sea $\mathfrak q\subseteq\mathbb Z[t]$ primo. Su contracción a $\mathbb Z$ es $(0)$ o $(p)$.

Si $\mathfrak q\cap\mathbb Z=(0)$, localizamos en $\mathbb Z\setminus\set{0}$. Por las \cref{prop:s08-primos-localizacion,prop:s08-polinomios-campo}, $\mathfrak q$ es la contracción de $(0)$ o de $(f)\subseteq\mathbb Q[t]$, con $f$ irreducible. En el primer caso es $(0)$. En el segundo, la \cref{prop:s08-gauss}(G1) permite sustituir $f$ por un primitivo irreducible $g\in\mathbb Z[t]$, único salvo signo. Por (G2),
\[
\mathfrak q=g\mathbb Q[t]\cap\mathbb Z[t]=g\mathbb Z[t].
\]
Estos ideales son primos y no contienen enteros no nulos.

Si $\mathfrak q\cap\mathbb Z=(p)$, el primo contiene a $(p)$. La \cref{prop:s08-primos-cociente} lo identifica con un primo de $\mathbb F_p[t]$: es $(0)$ o $(\bar g)$, con $\bar g$ mónico irreducible. Al contraer obtenemos $(p)$ o $(p,g)$ para cualquier levantamiento mónico $g$. Dos levantamientos difieren por un elemento de $p\mathbb Z[t]$, de modo que dan el mismo ideal. La unicidad del generador mónico en $\mathbb F_p[t]$ y las correspondencias usadas prueban que no hay más identificaciones.
\end{proof}

\begin{samepage}
\begin{prop}[Los puntos de $\mathbb P_{\mathbb Z}^1$]
\label{prop:s08-puntos-sobre-z}
\label{ejem:s08-puntos-sobre-z}
Los puntos de $\mathbb P_{\mathbb Z}^1$ son los primos de $U_0=\Spec\mathbb Z[t]$ de la \cref{prop:s08-primos-z-t} y los puntos de la sección del infinito $\psi_1(V(s))\cong\Spec\mathbb Z$: $(s)$ y los $(s,p)$, con $p$ primo. Sus cuerpos residuales y las imágenes por $\pi$ son:
\begin{center}
\normalfont
\small
\renewcommand{\arraystretch}{1.3}
\begin{tabular}{lccc}
\hline
Punto & Ideal & Cuerpo residual & Sobre\\
\hline
genérico & $(0)$ & $\mathbb Q(t)$ & $(0)$\\
cerrado en la fibra racional & $(g)$ & $\mathbb Q[t]/(g)$ & $(0)$\\
genérico de la fibra en $p$ & $(p)$ & $\mathbb F_p(t)$ & $(p)$\\
cerrado en la fibra en $p$ & $(p,g)$ & $\mathbb F_p[t]/(\bar g)$ & $(p)$\\
infinito, fibra racional & $(s)$ & $\mathbb Q$ & $(0)$\\
infinito, fibra en $p$ & $(s,p)$ & $\mathbb F_p$ & $(p)$\\
\hline
\end{tabular}
\end{center}
Las cuatro primeras filas son primos de $\mathbb Z[t]$ en $U_0$; las dos últimas son primos de $\mathbb Z[s]$ en $U_1$.
\end{prop}
\end{samepage}
\begin{proof}
La lista se obtiene de las \cref{prop:s08-primos-z-t,prop:s08-proyectiva-como-conjunto}. El cuerpo residual de un primo $\mathfrak q$ es $\Frac(B/\mathfrak q)$ en su carta $\Spec B$. Para $(g)$ se puede localizar primero en los enteros no nulos y obtener el campo $\mathbb Q[t]/(g)$. Para $(p)$ se obtiene $\mathbb F_p(t)$ y para $(p,g)$ el cociente $\mathbb F_p[t]/(\bar g)$, que ya es un campo. Los cocientes por $(s)$ y $(s,p)$ son $\mathbb Z$ y $\mathbb F_p$. Finalmente, el morfismo estructural contrae cada primo a $\mathbb Z$, lo que da la última columna.
\end{proof}

% Presentación S08, diapositiva 26: superficie aritmética y secciones.
\begin{center}
\begin{tikzpicture}[every node/.style={font=\small},x=1.05cm,y=.85cm]
\draw[grisCIMAT,->] (-.4,0)--(9.4,0) node[right] {$\Spec\mathbb Z$};
\foreach \x/\p in {2/2,4/3,6/5,8/7}{
 \draw[green!45!black,thick] (\x,.5)--(\x,3.7);
 \fill[grisCIMAT] (\x,0) circle (2pt) node[below] {$(\p)$};
}
\draw[blue!60!black,thick] plot[domain=.5:3.7,samples=75] ({.05*sin(700*\x)},\x);
\node[below] at (0,0) {$(0)$};
\node[above right,blue!60!black] at (0,3.7) {$\mathbb P_{\mathbb Q}^1$};
\node[above,green!45!black] at (2,3.7) {$\mathbb P_{\mathbb F_2}^1$};
\draw[orange!70!black,thick] (-.35,3.35)--(9,3.35) node[right] {$s=0$};
\draw[rojoCIMAT,thick] (-.35,1.35)--(9,1.35) node[right] {$t=0$};
\draw[orange!70!black,thick] (-.35,2.35) .. controls (.8,2.1) and (1.3,1.7) .. (2,1.35)
 .. controls (4,.6) and (6,.55) .. (9,.55) node[right] {$t=2$};
\fill[rojoCIMAT] (2,1.35) circle (3pt) node[above right] {$(2,t)$};
\node at (8.8,.65) {$\cdots$};
\draw[->,blue!60!black,thick] (-1,2.25)--(-1,.3) node[midway,left] {$\pi$};
\end{tikzpicture}
\end{center}
La base $\Spec\mathbb Z$ es una curva y cada fibra es una recta proyectiva: el dibujo representa $\mathbb P_{\mathbb Z}^1$ como superficie aritmética. Las secciones $t=0$ y $t=2$ se cortan solo sobre $(2)$, en el punto $(2,t)$.

\begin{prop}[Puntos cerrados de $\mathbb P_{\mathbb Z}^1$]
\label{prop:s08-puntos-cerrados-sobre-z}
\label{obs:s08-puntos-cerrados-sobre-z}
Un punto de $\mathbb P_{\mathbb Z}^1$ es cerrado si y solo si es de tipo $(p,g)$ o $(s,p)$, es decir, si cae sobre $(p)$ y es cerrado en la fibra $\mathbb P_{\mathbb F_p}^1$ de la \cref{prop:s08-fibras-proyectiva}. Su cuerpo residual es $\mathbb F_{p^d}$, donde $d$ es su grado en esa fibra. Para cada primo $p$ hay $p+1$ puntos cerrados con cuerpo residual $\mathbb F_p$ y $(p^2-p)/2$ con cuerpo residual $\mathbb F_{p^2}$.
\end{prop}
\begin{proof}
Los ideales $(p,g)$ son máximos porque sus cocientes son los campos finitos $\mathbb F_p[t]/(\bar g)$. Si el punto está en ambas cartas, el cambio $s=t^{-1}$ identifica también su cuerpo residual con un cociente de $\mathbb F_p[s]$; este es un campo, pues la inversa de un elemento no nulo de un campo finito es un polinomio en él. Así es cerrado en ambas cartas. El punto $(s,p)$ solo está en $U_1$ y es máximo allí. Aplicamos el \cref{lema:s08-cerrados-cartas}.

Los primos $(0)$ y $(p)$ no son máximos en $\mathbb Z[t]$. Tampoco lo es $(g)$: elijamos un primo $p$ que no divida al coeficiente principal de $g$. Entonces $\bar g$ es no constante y tiene un factor irreducible $\bar h\in\mathbb F_p[t]$. El ideal propio $(p,h)$ contiene estrictamente a $(g)$, pues $p\notin(g)$. Por último, $(s)$ no es máximo en $\mathbb Z[s]$. Esto excluye todos los demás puntos de la \cref{prop:s08-puntos-sobre-z}. Los grados y el conteo se deducen de las \cref{prop:s08-fibras-proyectiva,prop:s08-puntos-cerrados-campo-finito}.
\end{proof}

\begin{obs}[Dimensión]
\label{obs:s08-dimension-proyectiva-z}
La cadena de primos $(0)\subsetneq(p)\subsetneq(p,t)$ en $\mathbb Z[t]$ tiene longitud dos. Por tanto, la dimensión de Krull de $\mathbb Z[t]$ y la de $\mathbb P_{\mathbb Z}^1$, que contiene $\Spec\mathbb Z[t]$ como abierto, son al menos dos. Esto motiva la imagen de superficie; la dimensión se estudiará en la Sesión 16.
\end{obs}

\emph{Advertencia.} Un punto $(f)$ de la fibra racional es cerrado en $\mathbb P_{\mathbb Q}^1$, pero no en $\mathbb P_{\mathbb Z}^1$. Por ejemplo,
\[
\mathbb Z[t]/(2t-1)\cong\mathbb Z[1/2]
\]
no es un cuerpo. La cerradura de $(2t-1)$ es una curva horizontal que corta cada fibra una vez. En $U_0$ es $V(2t-1)$ y sus puntos corresponden, por la \cref{prop:s08-primos-cociente}, a los primos de $\mathbb Z[1/2]$: son el propio $(2t-1)$ y los $(p,2t-1)$ con $p$ impar. En la fibra sobre $2$ no hay punto en esta carta, porque $(2,2t-1)=(1)$. En $U_1$, la ecuación es $s=2$, de modo que la cerradura llega al infinito $(s,2)$.

\begin{samepage}
\subsection{Sobre un anillo de valoración discreta}

\begin{defi}[Anillo de valoración discreta]
\label{defi:s08-dvr}
Un anillo de valoración discreta (DVR) es un dominio de ideales principales local que no es un cuerpo. Sean $K$ su cuerpo de fracciones, $\mathfrak m=(\varpi)$ su único ideal máximo y $\kappa=A/\mathfrak m$ su cuerpo residual. El generador $\varpi$ se llama uniformizante.

Ejemplos son $A=\mathbb Z_{(p)}=\set{a/b\in\mathbb Q\given p\nmid b}$, con $K=\mathbb Q$, $\mathfrak m=(p)$ y $\kappa=\mathbb F_p$, y $A=k[u]_{(u)}$, con $K=k(u)$, $\varpi=u$ y $\kappa=k$. En $\mathbb Z_{(p)}$, las no unidades son los $a/b$ con $p\mid a$, así que su único ideal máximo es $p\mathbb Z_{(p)}$. Es un dominio de ideales principales por ser una localización de $\mathbb Z$.
\end{defi}
\end{samepage}

\begin{prop}[Estructura de un DVR]
\label{prop:s08-estructura-dvr}
Sea $A$ un DVR con uniformizante $\varpi$.
\begin{enumerate}[label=(\alph*)]
\item $\Spec A=\set{(0),\mathfrak m}$, $D(\varpi)=\set{(0)}$ y $\set{\mathfrak m}$ es cerrado.
\item Cada $x\in K^\times$ se escribe de manera única como $x=u\varpi^n$, con $u\in A^\times$ y $n\in\mathbb Z$. Definimos $v(x)=n$ y $v(0)=\infty$.
\end{enumerate}
\emph{[Atiyah--Macdonald, cap.~9; Sesión 3.]}
\end{prop}
\begin{proof}
Todo primo no nulo de un dominio de ideales principales es máximo, por la \cref{prop:s08-polinomios-campo}; como $A$ es local, debe ser $\mathfrak m$. Esto prueba (a). Además, un dominio de ideales principales tiene factorización única y todo irreducible de $A$ genera el único primo no nulo, por lo que es asociado a $\varpi$. Así, todo elemento no nulo de $A$ es una unidad por una potencia no negativa de $\varpi$. Escribiendo un elemento de $K^\times$ como cociente obtenemos (b). La unicidad de la factorización da la del exponente y, fijado este, la de la unidad.
\end{proof}

\begin{prop}[Las fibras de $\mathbb P_A^1$ sobre un DVR]
\label{prop:s08-fibras-dvr}
\label{ejem:s08-proyectiva-dvr}
Sea $A$ un DVR y $\pi:\mathbb P_A^1\to\Spec A$ el morfismo estructural.
\begin{enumerate}[label=(\alph*)]
\item La fibra genérica $\pi^{-1}((0))$ es abierta y, como conjunto, es $\mathbb P_K^1$.
\item La fibra especial $\pi^{-1}(\mathfrak m)$ es cerrada y, como conjunto, es $\mathbb P_\kappa^1$.
\item El punto $(0)$ de $A[t]$ es el genérico de $\mathbb P_A^1$, y la fibra genérica es densa.
\end{enumerate}
\end{prop}
\begin{proof}
La \cref{prop:s08-fibras-proyectiva} da las dos descripciones de puntos. La \cref{prop:s08-estructura-dvr}(a) y la continuidad de $\pi$ muestran que la primera fibra es abierta y la segunda cerrada. Como $A[t]$ es un dominio, $(0)$ es genérico en $U_0$; pertenece a $D(t)$ y se identifica con el genérico $(0)$ de $U_1$. Su cerradura contiene ambas cartas y es todo $\mathbb P_A^1$; como el punto pertenece a la fibra genérica, esta es densa.
\end{proof}

\begin{obs}[El disco]
\label{obs:s08-disco-dvr}
Se piensa $\Spec A$ como un disco pequeño. El punto cerrado $\mathfrak m$ es el centro, y el genérico $(0)$ representa al disco sin su centro. Así, $\mathbb P_A^1$ es una familia de rectas proyectivas sobre el disco: la fibra especial vive sobre el centro y la genérica sobre el resto.
\end{obs}

\begin{defi}[Puntos racionales y reducción]
\label{defi:s08-reduccion}
Para un campo $F$, $\mathbb P^1(F)$ es el conjunto de clases $[a:b]$ de pares $(a,b)\in F^2\setminus\set{(0,0)}$ módulo $F^\times$. El punto asociado de $\mathbb P_F^1$ es
\[
x_{[a:b]}=
\begin{cases}
(t-b/a),&a\ne0,\\
\infty,&a=0.
\end{cases}
\]
Si $b\ne0$, también tiene coordenada $s=a/b$ en $U_1$. Por la \cref{prop:s08-puntos-proyectiva-cuerpo}, esta aplicación es una biyección sobre los puntos de grado uno.

Sea ahora $A$ un DVR. Un representante $(a,b)$ de un punto de $\mathbb P^1(K)$ es normalizado si $a,b\in A$ y al menos uno es una unidad, es decir, si $\min(v(a),v(b))=0$. Su reducción es $[\bar a:\bar b]\in\mathbb P^1(\kappa)$, donde las barras indican reducción módulo $\mathfrak m$.
\end{defi}

\begin{lema}[La reducción está bien definida]
\label{lema:s08-reduccion-bien-definida}
\label{obs:s08-reduccion-bien-definida}
Todo punto de $\mathbb P^1(K)$ tiene un representante normalizado, y dos representantes normalizados difieren por una unidad de $A$. En consecuencia, $[\bar a:\bar b]$ es un punto de $\mathbb P^1(\kappa)$ independiente del representante.
\end{lema}
\begin{proof}
Sea $r=\min(v(a),v(b))$, que es finito porque las coordenadas no son ambas cero. Multiplicar por $\varpi^{-r}$ hace que ambas valoraciones sean no negativas y una sea cero. Por la \cref{prop:s08-estructura-dvr}(b), el nuevo par está normalizado.

Si $(a',b')=\lambda(a,b)$ y ambos pares están normalizados, entonces
\[
0=\min(v(a'),v(b'))=v(\lambda)+\min(v(a),v(b))=v(\lambda).
\]
Así $\lambda\in A^\times$ y $\bar\lambda\in\kappa^\times$. Los pares reducidos son no nulos y difieren por $\bar\lambda$, por lo que representan el mismo punto.
\end{proof}

Para $A=\mathbb Z_{(p)}$, esta es la reducción de un punto racional módulo $p$. Que todo punto $K$-racional se extienda a un punto con valores en $A$ es la parte de existencia del criterio valuativo de propiedad de la Sesión 15.

% S08b, diapositiva 33 y esqueleto, sección 5.3: extensión sobre el DVR.
\begin{center}
\begin{tikzcd}[column sep=large,row sep=large]
\Spec K\arrow[r,"{[a:b]}"]\arrow[d]
 &\mathbb P_A^1\arrow[d,"\pi"]\\
\Spec A\arrow[r,swap,"\id"]\arrow[ur,dashed]&\Spec A
\end{tikzcd}
\end{center}
La flecha discontinua representa la extensión a $\Spec A$ del punto definido sobre $K$; su valor en el punto cerrado es la reducción.

\begin{prop}[La cerradura corta la fibra especial en la reducción]
\label{prop:s08-cerradura-reduccion}
Sea $[a:b]\in\mathbb P^1(K)$ y $x=x_{[a:b]}$ el punto correspondiente de la fibra genérica de $\mathbb P_A^1$. La cerradura de $\set{x}$ corta la fibra especial en exactamente un punto: el que la \cref{prop:s08-fibras-proyectiva} identifica con $x_{[\bar a:\bar b]}\in\mathbb P_\kappa^1$, para un representante normalizado $(a,b)$.
\end{prop}
\begin{proof}
Normalizamos las coordenadas mediante el \cref{lema:s08-reduccion-bien-definida}. Supongamos primero que $a$ es una unidad y pongamos $c=b/a\in A$. En $U_0$, el punto $x$ es el primo $(t-c)$ de $A[t]$: el cociente es $A$ y su contracción a $A$ es $(0)$. Su cerradura en $U_0$ es $V(t-c)\cong\Spec A$ y corta la fibra especial en el único primo
\[
\mathfrak mA[t]+(t-c),
\]
que corresponde al punto $t=\bar c=\bar b/\bar a$.

No aparece otro punto en la segunda carta. Si $b=0$, entonces $x\notin U_1$ y la \cref{obs:s08-especializacion-local} muestra que su cerradura tampoco corta a $U_1$. Si $b\ne0$, entonces $x\in U_1$ y la ecuación $bs-a=0$ se cumple en $x$, luego también en su cerradura dentro de esa carta. Cuando $b\in\mathfrak m$, su reducción es $-\bar a=0$, imposible, así que no hay punto de la fibra especial en $U_1$. Si $b$ es unidad, la cerradura es $V(s-a/b)$; su punto especial $s=\bar a/\bar b$ es el ya encontrado en el traslape.

Si $a$ no es unidad, $b$ sí lo es. El mismo argumento, intercambiando las cartas, da un único punto especial con $s=\bar a/\bar b$. En todos los casos es el punto $[\bar a:\bar b]$.
\end{proof}

% Presentación S08, diapositiva 30: especialización de un punto racional.
\begin{center}
\begin{minipage}{.45\textwidth}
\centering
\begin{tikzpicture}[every node/.style={font=\small}]
\draw[grisCIMAT] (-.7,0)--(3.3,0);
\draw[blue!60!black,thick] plot[domain=.35:2.9,samples=65] ({.05*sin(700*\x)},\x);
\draw[green!45!black,thick] (2.3,.35)--(2.3,2.9);
\node[above,blue!60!black] at (0,2.9) {$\mathbb P_K^1$};
\node[above,green!45!black] at (2.3,2.9) {$\mathbb P_\kappa^1$};
\fill[blue!60!black] (0,0) circle (2pt) node[below] {$(0)$};
\fill[green!45!black] (2.3,0) circle (2pt) node[below] {$\mathfrak m$};
\draw[rojoCIMAT,thick] (0,1.85) .. controls (.8,2.2) and (1.7,2.15) .. (2.3,1.6);
\fill[rojoCIMAT] (0,1.85) circle (2.5pt) node[left] {$[a:b]$};
\fill[rojoCIMAT] (2.3,1.6) circle (2.5pt) node[right] {$[\bar a:\bar b]$};
\end{tikzpicture}
\end{minipage}\hfill
\begin{minipage}{.51\textwidth}
\small
La cerradura de un punto $K$-racional corta la fibra especial en su reducción. Si $A=k[u]_{(u)}$, se interpreta como el límite cuando $u\to0$. Para $A=\mathbb Z_{(p)}$, con $p$ impar, $1/2=[2:1]$ se reduce a $t=2^{-1}$ en $\mathbb F_p$.
\end{minipage}
\end{center}
\emph{Advertencia.} Para $A=\mathbb Z_{(2)}$, el punto $1/2=[2:1]$ se reduce a $[0:1]=\infty$. Hay que usar un representante con una coordenada unidad: $[1:1/2]$ no sirve.

\section{\texorpdfstring{$\mathbb P_A^n$}{PnA} por cartas}

\begin{defi}[Las cartas de $\mathbb P_A^n$]
\label{defi:s08-cartas-proyectiva-n}
Sea $n\geq1$. Para $i=0,\ldots,n$, sea $A_i=A[x_{k/i}:k\ne i]$, el anillo de polinomios en las $n$ variables $x_{k/i}$, y $U_i=\Spec A_i$. Se pone $x_{i/i}=1$, $U_{ii}=U_i$ y $\varphi_{ii}=\id$. Para $j\ne i$, sea $U_{ij}=D(x_{j/i})\subseteq U_i$, con anillo de funciones $A_i[x_{j/i}^{-1}]$, y
\[
\varphi_{ij}^\sharp:A_j[x_{i/j}^{-1}]\longrightarrow A_i[x_{j/i}^{-1}],
\qquad x_{k/j}\longmapsto\frac{x_{k/i}}{x_{j/i}}.
\]
Es un isomorfismo: envía $x_{i/j}$ a $x_{j/i}^{-1}$, con inversa $x_{k/i}\mapsto x_{k/j}/x_{i/j}$. \emph{[Vakil, \S4.4.9 y \S4.4.10.]}
\end{defi}

\begin{prop}
\label{prop:s08-cociclo-proyectiva-n}
Los $U_i$, $U_{ij}$ y $\varphi_{ij}$ son datos de pegado, incluida la igualdad de dominios de (P3):
\[
\varphi_{ij}(U_{ij}\cap U_{ik})=U_{ji}\cap U_{jk}.
\]
El esquema del \cref{teo:s08-lema-pegado} es el espacio proyectivo $\mathbb P_A^n$.
\end{prop}

\emph{Advertencia.} (P3) exige dos cosas: que $\varphi_{ij}$ lleve exactamente $U_{ij}\cap U_{ik}$ sobre $U_{ji}\cap U_{jk}$, y que el triángulo de transiciones conmute. La primera no se sigue de la segunda.

% Presentación S08, diapositiva 32: dominios y triángulo del cociclo.
\begin{center}
\begin{tikzcd}[column sep=small,row sep=large]
& U_{ij}\cap U_{ik}\subseteq U_i
  \arrow[dl,swap,"\varphi_{ij}"] \arrow[dr,"\varphi_{ik}"] &\\
U_{ji}\cap U_{jk}\subseteq U_j \arrow[rr,swap,"\varphi_{jk}"]
&& U_{ki}\cap U_{kj}\subseteq U_k
\end{tikzcd}
\end{center}
\begin{proof}
Sea
\[
L=A[x_0^{\pm1},\ldots,x_n^{\pm1}],\qquad
L_0=\bigoplus_{\substack{m\in\mathbb Z^{n+1}\\ \sum_\ell m_\ell=0}}Ax^m,
\quad x^m=x_0^{m_0}\cdots x_n^{m_n}.
\]
El submódulo $L_0$ es un subanillo. Para cada $i$, definimos
\[
\iota_i:A_i\hookrightarrow L_0,\qquad x_{k/i}\longmapsto x_kx_i^{-1}.
\]
Es inyectivo porque monomios distintos se envían a monomios distintos. La imagen de $A_i[x_{j/i}^{-1}]$ es
\[
R_{ij}=\bigoplus_{\substack{\sum_\ell m_\ell=0\\m_\ell\geq0\text{ si }\ell\notin\set{i,j}}}Ax^m.
\]
En efecto, los exponentes fuera de $i,j$ son no negativos; el exponente en $j$ puede ser cualquiera, y el grado total cero fija el exponente en $i$. La descripción es simétrica en $i,j$. Además,
\[
\iota_j(x_{k/j})=\iota_i(x_{k/i}/x_{j/i}),
\qquad \varphi_{ij}^\sharp=\iota_i^{-1}\iota_j
\]
en las localizaciones identificadas con $R_{ij}$. Así se verifican (P1) y (P2).

El anillo del traslape triple en la carta $i$ se identifica con
\[
R_{ijk}=\bigoplus_{\substack{\sum_\ell m_\ell=0\\m_\ell\geq0\text{ si }\ell\notin\set{i,j,k}}}Ax^m,
\]
al invertir $x_{j/i}$ y $x_{k/i}$. En la carta $j$, los abiertos $U_{ji}\cap U_{jk}$ tienen el mismo anillo: bajo $\varphi_{ij}^\sharp$, invertir $x_{k/j}$ equivale a invertir $x_{k/i}/x_{j/i}$, y $x_{j/i}$ ya es unidad. Por tanto la transición lleva exactamente el primer abierto triple al segundo. Las tres transiciones inducen la identidad de $R_{ijk}$; su composición coincide con la transición directa. Esto prueba ambas partes de (P3).
\end{proof}

\begin{obs}
\label{obs:s08-hiperplano-infinito}
Geométricamente, $U_i$ es el abierto donde $x_i\ne0$. La carta $U_0$ es isomorfa a $\mathbb A_A^n$. Su complemento es el cerrado donde se anula $x_0$, descrito en $U_i$ por $V(x_{0/i})$. Es una copia de $\mathbb P_A^{n-1}$, el hiperplano al infinito de la Sesión 10. Para $n=1$ se recupera la \cref{defi:s08-recta-proyectiva}, con $t=x_{1/0}$ y $s=x_{0/1}$. La construcción formal mediante $\operatorname{Proj}$ llegará en la Sesión 21.
\end{obs}

\section{Los pegados de dos rectas a lo largo de \texorpdfstring{$D(x)$}{D(x)}}
\label{sec:s08-dos-rectas}

\begin{lema}[Funtorialidad del pegado]
\label{lema:s08-funtorialidad-pegado}
Sean $(X_i,X_{ij},\varphi_{ij})$ y $(X_i',X_{ij}',\varphi_{ij}')$ datos de pegado con los mismos índices, y pegados $(X,\psi_i)$ y $(X',\psi_i')$. Sean $\nu_i:X_i\to X_i'$ isomorfismos tales que $\nu_i(X_{ij})=X_{ij}'$ y $\nu_j\circ\varphi_{ij}=\varphi_{ij}'\circ\nu_i$ en $X_{ij}$. Existe un único isomorfismo $h:X\to X'$ con $h\circ\psi_i=\psi_i'\circ\nu_i$.
\end{lema}

% Presentación S08, diapositiva 34: los dos cuadrados de funtorialidad.
\begin{center}
\begin{minipage}{.43\textwidth}
\centering
\begin{tikzcd}
X_i\arrow[r,"\psi_i"]\arrow[d,swap,"\nu_i"]
 &X\arrow[d,"h"]\\
X_i'\arrow[r,swap,"\psi_i'"]&X'
\end{tikzcd}
\par\smallskip{\small La conclusión, para todo $i$.}
\end{minipage}\hfill
\begin{minipage}{.53\textwidth}
\centering
\begin{tikzcd}
X_{ij}\arrow[r,"\varphi_{ij}"]\arrow[d,swap,"\nu_i"]
 &X_{ji}\arrow[d,"\nu_j"]\\
X_{ij}'\arrow[r,swap,"\varphi_{ij}'"]&X_{ji}'
\end{tikzcd}
\par\smallskip{\small La hipótesis.}
\end{minipage}
\end{center}
\begin{proof}
Los morfismos $\psi_i'\circ\nu_i:X_i\to X'$ son compatibles, pues
\[
\psi_j'\circ\nu_j\circ\varphi_{ij}
=\psi_j'\circ\varphi_{ij}'\circ\nu_i
=\psi_i'\circ\nu_i.
\]
La \cref{obs:s08-morfismos-desde-pegado} da un único $h:X\to X'$ con la propiedad requerida. Los $\nu_i^{-1}$ satisfacen las mismas condiciones en sentido inverso y producen $h':X'\to X$. Sus composiciones restringen a la identidad en cada carta; como las cartas cubren, $h'h=\id_X$ y $hh'=\id_{X'}$.
\end{proof}

\begin{prop}[Pegados de dos rectas]
\label{prop:s08-pegados-dos-rectas}
Sea $k$ un campo, $X_1=\Spec k[x]$, $X_2=\Spec k[y]$, $X_{12}=D(x)$ y $X_{21}=D(y)$. Para un isomorfismo de $k$-esquemas $\varphi:X_{12}\to X_{21}$, el pegado es isomorfo a la recta con origen doble si $\varphi^\sharp(y)=cx$, o a $\mathbb P_k^1$ si $\varphi^\sharp(y)=cx^{-1}$, con $c\in k^\times$. No hay otros casos, y los dos esquemas no son isomorfos.
\end{prop}
\begin{proof}
El isomorfismo de $k$-álgebras
\[
\varphi^\sharp:k[y,y^{-1}]\xrightarrow{\sim}k[x,x^{-1}]
\]
envía $y$ a una unidad. Las unidades de $k[x,x^{-1}]$ son los monomios $cx^n$, con $c\in k^\times$ y $n\in\mathbb Z$. En un producto de polinomios de Laurent sobre un campo, los exponentes mínimo y máximo se suman. Si el producto es uno, cada factor tiene un solo exponente.

La inversa también envía $x$ a una unidad, digamos $dy^m$. Al componer,
\[
x=dc^m x^{nm},
\]
de modo que $nm=1$ y $n=\pm1$. Tomemos $\nu_1=\id$ y el cambio de coordenada en la segunda carta definido por $\nu_2^\sharp(y')=c^{-1}y$. Entonces
\[
\varphi^\sharp\nu_2^\sharp(y')=c^{-1}cx^{\pm1}=x^{\pm1}.
\]
Además, $\nu_2^\sharp$ lleva $(y')^{\pm1}$ a $c^{\mp1}y^{\pm1}$, de modo que $\nu_2(X_{21})=X_{21}'$. La compatibilidad para el par inverso se obtiene invirtiendo la del par $(1,2)$. El \cref{lema:s08-funtorialidad-pegado} identifica el pegado con el de $y'=x$ o el de $y'=x^{-1}$. Son la recta con origen doble del \cref{ejem:s07-recta-origen-doble} y la recta proyectiva de la \cref{defi:s08-recta-proyectiva}. La constante se absorbe en el cambio de carta; el exponente distingue los casos. Que no son isomorfos se demuestra en la \cref{prop:s08-secciones-origen-doble}.
\end{proof}

\begin{prop}[Secciones globales de la recta con origen doble]
\label{prop:s08-secciones-origen-doble}
Sea $X$ la recta con origen doble sobre un campo $k$, con cartas $X_1=\Spec k[x]$ y $X_2=\Spec k[y']$ pegadas por $y'=x$ en los abiertos donde la coordenada es invertible. Entonces $\Gamma(X,\mathcal O_X)=k[x]$. En consecuencia, $X$ y $\mathbb P_k^1$ no son isomorfos, y $X$ no es afín. \emph{[Vakil, ejercicio 4.4.B.]}
\end{prop}
\begin{proof}
Una sección es un par $(f(x),g(y'))$ que coincide en $k[x,x^{-1}]$ tras sustituir $y'=x$. Como $k[x]\hookrightarrow k[x,x^{-1}]$ es inyectivo, la condición equivale a $f=g$. Así, las secciones son los pares $(f,f)$ y su anillo es $k[x]$. Un isomorfismo con $\mathbb P_k^1$ induciría un isomorfismo entre $k[x]$ y $k$, por la \cref{prop:s08-secciones-proyectiva}; esto es imposible, pues $k$ es un campo y $k[x]$ no.

Las restricciones $\Gamma(X,\mathcal O_X)\to\Gamma(X_i,\mathcal O_{X_i})$ son isomorfismos. Por la construcción de la \cref{defi:s08-morfismo-canonico}, $\eta_X:X\to\Spec k[x]$ restringe a la identidad en cada carta, identificando $y'$ con $x$. Envía los dos orígenes distintos al mismo punto $(x)$, de modo que no es inyectivo. La \cref{prop:s08-propiedades-morfismo-canonico}(e) implica que $X$ no es afín.
\end{proof}

\begin{obs}[Los dos pegados sobre un anillo]
\label{obs:s08-dos-pegados-anillo}
El mismo cálculo de secciones de la recta con origen doble vale sobre cualquier anillo $A$: $A[x]\hookrightarrow A[x,x^{-1}]$ sigue siendo inyectivo y el anillo de secciones es $A[x]$. Si $A\ne0$, para cualquier $\mathfrak p\in\Spec A$ los dos puntos $\mathfrak pA[x]+(x)$, uno en cada carta, son distintos y tienen la misma imagen por el morfismo canónico. Por tanto, la recta con origen doble tampoco es afín sobre $A$. Junto con el \cref{coro:s08-proyectiva-no-afin}, esto da dos pegados no afines de las mismas piezas. Sobre un campo los distingue la \cref{prop:s08-secciones-origen-doble}; para un anillo no nulo cualquiera los distinguirá la separación: $\mathbb P_A^1$ es separado y la recta con origen doble no lo es. Esta propiedad se estudiará en la Sesión 15.
\end{obs}

\section{Ejercicios}

\begin{ejer}[El espacio afín menos el origen]
\label{ejer:s08-afin-menos-origen}
Sea $k$ un campo, $n\geq2$ y $U=\mathbb A_k^n\setminus\set{(x_1,\ldots,x_n)}$. Muestra que $\Gamma(U,\mathcal O_U)=k[x_1,\ldots,x_n]$ y que $U$ no es afín. Pista: cubre $U$ con $D(x_1),\ldots,D(x_n)$ y adapta la \cref{prop:s08-secciones-plano-menos-origen}. \emph{[Cf. Vakil, \S4.4.1 a \S4.4.3 (caso $n=2$) y \S4.4.2.]}
\end{ejer}

\begin{ejer}[La recta menos el origen]
\label{ejer:s08-recta-menos-origen}
Sea $k$ un campo y $V=\mathbb A_k^1\setminus\set{0}=D(x)$. Calcula $\Gamma(V,\mathcal O_V)$ y describe $\eta_V$. Justifica por qué es un isomorfismo, a diferencia del plano menos el origen.
\end{ejer}

\begin{ejer}[El automorfismo que intercambia las cartas]
\label{ejer:s08-intercambio-cartas}
Muestra que $t\leftrightarrow s$ define un automorfismo de $\mathbb P_A^1$ que intercambia $\psi_0(U_0)$ y $\psi_1(U_1)$. Pista: usa el \cref{lema:s08-funtorialidad-pegado} con $\nu_0:U_0\to U_1$, $\nu_1:U_1\to U_0$ y, como segundos datos de pegado, los mismos con los índices $0$ y $1$ intercambiados.
\end{ejer}

\begin{ejer}[Puntos de $\mathbb P_{\mathbb F_2}^1$ y $\mathbb P_{\mathbb F_3}^1$]
\label{ejer:s08-puntos-f2-f3}
Lista los puntos cerrados de $\mathbb P_{\mathbb F_q}^1$ de grado uno y dos para $q=2$ y $q=3$. Verifica $M_2=(q^2-q)/2$ y encuentra $M_3$ para $q=2$.
\end{ejer}

\begin{ejer}[Los puntos cerrados de $\mathbb P_{\mathbb Z}^1$ sobre $2$]
\label{ejer:s08-puntos-sobre-dos}
Describe la fibra de $\mathbb P_{\mathbb Z}^1\to\Spec\mathbb Z$ sobre $(2)$ y lista sus puntos cerrados con cuerpo residual $\mathbb F_2$ y $\mathbb F_4$.
\end{ejer}

\begin{ejer}[$\mathbb P_A^n$]
\label{ejer:s08-proyectiva-n}
Muestra que $\Gamma(\mathbb P_A^n,\mathcal O)=A$ y que $\mathbb P_A^n$ no es afín si $n\geq1$ y $A\ne0$. Pista: una sección global se restringe de forma inyectiva a $U_0$, pues $A_i\to A_i[x_{0/i}^{-1}]$ es inyectivo. Basta entonces intersecar los anillos de $U_0$ y $U_1$ dentro de $L_0$, como en la \cref{prop:s08-cociclo-proyectiva-n}. \emph{[Vakil, ejercicio 4.4.E.]}
\end{ejer}
